% =====================================================================
%  Limits of Metamorphic Testing for Object-Level Authorization:
%  Temporal Reachability, Oracle Admissibility, and Measurement Resolution
%  Official IEEE Access template (ieeeaccess.cls, which \LoadClass{IEEEtran})
% ---------------------------------------------------------------------
%  硬规则:
%    * 正文**不得手敲任何数字** —— 全部引用 \Res* 宏（由
%      04_绘图与分析/analysis/compute_all.py 从原始记录重算生成）。
%    * 参考文献**不得手敲** —— refs.bib 由
%      04_绘图与分析/analysis/normalize_bib.py 从带 DOI 的记录规范化生成。
%    * 全部图为**矢量**，由 matplotlib / figure-spec 确定性生成，无 AI 生成像素。
%  模板: paper/ieeeaccess/ 是官方 IEEE Access LaTeX 模板（ieeeaccess.cls +
%        IEEEtran.cls/.bst + Formata/Times 字体 + 位图），来源
%        ACCESS_latex_template_20260513-1-1.zip（IEEE Access 官方下载）。
%        该类**不会**自动加载 hyperref（仅 \ifjtehm 分支加载），故此处显式加载，
%        以支持 \orcidicon（ORCID iD 图标）与 PDF 元数据。
%  编译: bash build.sh   （会把源码与模板支持文件装配到纯 ASCII 目录再编译）
% =====================================================================
\documentclass{ieeeaccess}

\usepackage{cite}
\usepackage{amsmath}
\usepackage{amssymb}
\usepackage{amsthm}
\usepackage{booktabs}
\usepackage{graphicx}
\usepackage{url}
\usepackage{etoolbox}
\usepackage[hidelinks]{hyperref}
% ORCID iD 图标：ieeeaccess.cls 把 TeX 原始量 \year 重定义成带参宏
%   （\year#1->\gdef\theyear{#1}），而 pgf/pgfmath 加载时要读 \year 取随机种子，
%   于是 tikz/orcidlink 会直接 Fatal error —— 二者**不可共存**。
%   这里改用 pict2e 直接绘制绿色 iD 徽标（纯 LaTeX 代码，无外部文件、字节可复现）。
\usepackage{pict2e}
\definecolor{orcidgreen}{rgb}{0.65098,0.80784,0.22353}% ORCID iD 官方绿 A6CE39
\newsavebox{\ORCIDlogo}
\savebox{\ORCIDlogo}{%
  \setlength{\unitlength}{\dimexpr 1em/256\relax}%
  \begin{picture}(256,256)%
    \color{orcidgreen}\put(128,128){\circle*{256}}%
    \color{white}%
    \put(78.6,199.2){\circle*{20}}%
    \moveto(70.9,176.9)\lineto(86.3,176.9)\lineto(86.3,69.8)\lineto(70.9,69.8)%
    \closepath\fillpath%
    \moveto(108.9,176.9)\lineto(150.5,176.9)%
    \curveto(190.1,176.9)(207.5,148.6)(207.5,123.3)%
    \curveto(207.5,95.8)(186,69.7)(150.7,69.7)%
    \lineto(108.9,69.7)\closepath\fillpath%
    \color{orcidgreen}%
    \moveto(124.3,83.6)\lineto(148.8,83.6)%
    \curveto(183.7,83.6)(191.7,110.1)(191.7,123.3)%
    \curveto(191.7,144.8)(178,163)(148,163)%
    \lineto(124.3,163)\closepath\fillpath%
  \end{picture}%
}
% 用法：\orcidicon{0009-0000-5392-4475} —— 在作者名右上渲染可点击的 iD 徽标。
\newcommand{\orcidicon}[1]{\href{https://orcid.org/#1}{\usebox{\ORCIDlogo}}}
\hypersetup{
  pdftitle={Two Reachability-Gated Relations in MST-wi: Temporal Mismatch, Oracle Admissibility, and Measurement Limits},
  pdfauthor={Qiulong Lin and Nanlei Chen},
  pdfsubject={Metamorphic testing, object-level authorization, API security}
}

% ---- 图中/文中引用的全部数字（自动生成，勿手改） ----
\input{numbers.tex}
% ---- 原生引擎运行的数字（由 analysis/ingest_native_engine.py 生成） ----
\input{native_engine_numbers.tex}
% ---- GitLab 跨厂商复制的数字（由 analysis/ingest_gitlab.py 生成） ----
\input{gitlab_numbers.tex}
% ---- Gogs 跨厂商复制的数字与对照表（由 analysis/ingest_gogs.py 生成） ----
\input{gogs_numbers.tex}
% ---- Gogs 按渠道/按机制的 oracle 保真矩阵（由 analysis/ingest_gogs_channel.py 生成） ----
\input{gogs_channel_numbers.tex}
% ---- 受控靶机阳性对照的数字与对照表（由 analysis/ingest_seeded.py 生成） ----
\input{seeded_table.tex}
% ---- 真实 vulnerable/fixed 配对（由 analysis/ingest_cve_pair.py 生成） ----
\input{cve_pair_numbers.tex}
% ---- 真时序转移实验（由 analysis/ingest_temporal.py 生成） ----
\input{temporal_numbers.tex}
% ---- M2 补充分析：oracle 有效性分层（由 analysis/ingest_m2_stratum.py 生成） ----
\input{m2_stratum.tex}
% ---- 两条关系的显式谓词与真值表（由 analysis/ingest_truth.py 生成） ----
\input{truth_numbers.tex}
% ---- 完整因子设计、配对归因与 36/18 对账（由 analysis/ingest_design.py 生成） ----
\input{design_numbers.tex}
% ---- 覆盖率/弃权/选择性风险/成本（由 analysis/ingest_coverage.py 生成） ----
\input{coverage_numbers.tex}
% ---- 竞争基线：常规多账户差分（由 analysis/ingest_baseline.py 生成） ----
\input{baseline_numbers.tex}
% ---- 等总预算比较与多类缺陷（由 ingest_budget.py / ingest_multiclass.py 生成） ----
% 摘要引用了 \McInstances，故 multiclass 宏必须在正文之前加载；
% budget 宏只被 budget_table.tex 使用，由该表自行加载，这里不重复。
\input{multiclass_numbers.tex}

% 图形搜索路径：figures/ 放正文插图；./ 供 ieeeaccess.cls 内部取 bullet.png
% （摘要/关键词前的圆点、\EOD 结尾标记）。若只写 {figures/} 会让类找不到 ./bullet.png。
\graphicspath{{figures/}{./}}

% 双栏窄列下允许更宽的词间弹性，避免长 token 造成 overfull hbox
\setlength{\emergencystretch}{3em}

% ---- 版面微调（控制页数，不改字号、不改行距、不改栏宽） ----
% 只压缩浮动体/公式/表格周围的**垂直留白**，不影响可读性，也不改变任何正文内容。
\setlength{\textfloatsep}{6pt plus 2pt minus 2pt}
\setlength{\floatsep}{6pt plus 2pt minus 2pt}
\setlength{\intextsep}{6pt plus 2pt minus 2pt}
\setlength{\abovecaptionskip}{3pt plus 1pt minus 1pt}
\setlength{\belowcaptionskip}{0pt}
\setlength{\abovedisplayskip}{4pt plus 1pt minus 1pt}
\setlength{\belowdisplayskip}{4pt plus 1pt minus 1pt}
\setlength{\abovedisplayshortskip}{2pt plus 1pt}
\setlength{\belowdisplayshortskip}{2pt plus 1pt}
\renewcommand{\arraystretch}{0.95}

% 表格与图题降为 \footnotesize，压缩浮动体高度
\AtBeginEnvironment{table}{\footnotesize}
\AtBeginEnvironment{table*}{\footnotesize}
\AtBeginEnvironment{figure}{\footnotesize}
\AtBeginEnvironment{figure*}{\footnotesize}

% 可复现构建：省略 PDF 的 CreationDate/ModDate 以及嵌入的本地路径信息，
% 使同一份源码每次编译得到字节一致的 main.pdf（这些原语仅在 pdfTeX 下存在）。
\ifdefined\pdfinfoomitdate\pdfinfoomitdate=1\fi
\ifdefined\pdfsuppressptexinfo\pdfsuppressptexinfo=-1\fi

\newcommand{\ttt}[1]{\texttt{#1}}
\newcommand{\code}[1]{\texttt{#1}}

% 形式化结果（引理/命题/推论）
\theoremstyle{definition}
\newtheorem{lemma}{Lemma}
\newtheorem{proposition}{Proposition}
\newtheorem{corollary}{Corollary}

% ---- 修复 ieeeaccess.cls 的 \paragraph 编号 ----------------------------------
% 该类把 \paragraph 当作**带字母编号**的第四级标题（\@seccntformat 取
% \theparagraphdis = \alph{paragraph}:），于是 "\paragraph{Ground truth}"
% 会渲染成 "a: Ground truth"。本文用它作无编号的行内引导短语，故去掉编号，
% 恢复为与正文同级的斜体行内标题；\section/\subsection/\subsubsection 不受影响。
\makeatletter
\renewcommand\paragraph{\@startsection{paragraph}{4}{\z@}%
  {-1.1ex\@plus-0.4ex\@minus-0.2ex}{0.35ex}{\normalfont\normalsize\itshape}}
\def\@seccntformat#1{%
  \ifstrequal{#1}{paragraph}{}{\csname the#1dis\endcsname\hskip 0.5em\relax}}
\makeatother

\begin{document}

% --- IEEE Access 首页元信息 -------------------------------------------------
% \history（出版/版本日期）与 DOI 由 IEEE 在录用排版阶段填写。此处**留空**，
% 不编造 DOI。注意：ieeeaccess.cls 的 \@maketitle 无条件使用 \@doi，
% 若不调用 \doi{} 会触发 "Undefined control sequence"，故必须给出（空值）。
\history{}
\doi{}

\title{Two Reachability-Gated Relations in MST-wi: Temporal Mismatch, Oracle Admissibility, and Measurement Limits}

% 作者块：IEEE Access 规范为「姓名大写 + 上标单位编号」；ORCID 以 iD 图标贴在
% 姓名右上。Lin 的 ORCID iD 由作者本人提供（2026-09-28）：0009-0000-5392-4475。
% Chen 尚未提供 ORCID，故此处只挂 Lin 一个，不臆造。
\author{\uppercase{Qiulong Lin}\orcidicon{0009-0000-5392-4475}\authorrefmark{1,2}
AND \uppercase{Nanlei Chen}\authorrefmark{1,2}}

\address[1]{Yango University, Fuzhou 350015, China (e-mail: ubfbzu3@163.com)}
\address[2]{Sogang University, Seoul 04107, Republic of Korea}

\markboth{IEEE Access}{Lin and Chen: Reachability-Gated Relations in MST-wi}

\corresp{Corresponding author: Nanlei Chen (e-mail: ubfbzu3@163.com).}

\begin{abstract}
Two system-agnostic relations in MST-wi, a Web-system metamorphic testing
framework, gate their verdict on a \emph{GUI-reachability model}---a crawl-time
snapshot of what the interface permits---so what they establish about
object-level authorization depends on when the model was read and what the
adjudicator may consult. Source inspection shows the engine records only that a
relation failed: promotion to a \emph{proven violation} is added policy, stated
here explicitly. On a \ResSynRecords{}-record constructed study, a legitimate
access whose grant lies outside the framework's record is reported as proven in
\ResObsOobVtwoFalseProof{} of \ResObsOobFalseAlarm{} alarms, and the same
failure appears when the grant is made between crawl and request. Once a
readable, sound oracle is supplied, the relation adds no observed adjudication
benefit over the direct consumer of the same oracle and outcome; and the
published specificity of \ResLitReportedSpecPct{}\% is an input-level figure
whose denominator cannot, on its own, expose a class-uniform failure. A
temporal experiment on Gitea \ResRsVersion{} and GitLab CE \GlVersion{} commits
\TmpBothPolicyEvents{} authorisation transitions over \TmpBothCells{}
cross-credential cells: \TmpBothStaleFP{} stale false proofs appear, and a
recrawl with the same credentials removes \TmpBothFixed{}. A real
vulnerable/fixed pair (Gitea \CveVulnVersion{} to \CveFixedVersion{}) shows the
independent entitlement/effect gate reporting \CveTruePositives{} true positive
on the vulnerable build, \CveFalseProofs{} false proofs, and no alarm on the
fixed build---a check of the gate, not a measure of detection effectiveness;
across four seeded defect classes (\McInstances{} instances) the studied relation
fires only on the read class, a boundary rather than a tuning gap.
\end{abstract}

\begin{keywords}
BOLA/BFLA, false proof, metamorphic testing, object-level authorization, oracle problem, specificity measurement
\end{keywords}

\titlepgskip=-21pt

\maketitle

\section{Introduction}
\label{sec:intro}

\PARstart{B}{roken} object-level authorization (BOLA), first in the OWASP API
Security Top~10~\cite{OWASP2023API1}, is the failure to enforce, per object, that the
authenticated subject is entitled to the object touched. Detecting it is an
\emph{oracle} problem: vulnerable and correct systems answer with the same
success status. Metamorphic testing is an established approach to that oracle
problem, and \mbox{MST-wi}~\cite{Chaleshtari2022Metamorphic} is a substantial
system-agnostic instance for Web systems: \ResLitMrCatalog{} system-agnostic
relations (MRs) over \ResLitSmrlFunctions{} Web-specific predicates, reported to
detect \ResLitReportedDetection{}\% of vulnerabilities with specificity
\ResLitReportedSpecPct{}\%.

This paper asks what happens to that claim when the asserted property is
object-level authorization. The predicates do not express entitlement; they
express what their authors call reachability---what ``the functionality exposed
by the user interface'' permits---with the role-to-resource mapping
``\textbf{not always available}''~\cite{Chaleshtari2022Metamorphic}. We take
that sentence as premise, not discovery. Delegating adjudication to a reachability
model makes the moment the model was read part of the specification, so we
instantiate the predicates \emph{as defined}---over a model recorded at crawl
time---and drive the real systems through three phases: crawl, change
authorization out of band, execute. The middle phase is the one the original
evaluation does not exercise. The constructed study isolates the same defect
without time: it seeds a share outside the framework's record and asks whether
a gate that cannot see it still calls the access legal.

We ran a \ResSynRecords{}-record scenario/configuration study on a constructed
system and a controlled study on Gitea \ResRsVersion{}, extended, with a
methodological variation, to Gogs
\GgVersion{} and GitLab \GlVersion{}. We propose no mechanism. What survives is
three findings, one per research question (Section~\ref{sec:findings}).

\paragraph{Research questions and findings}
(1)~\emph{Temporal (RQ1).} Does the crawl-time reachability model remain sound
for authorization decisions after an out-of-band authorization change? No: the
adjudication as designed produces \emph{deterministic false proofs}---a legitimate
access the model cannot see is reported as a proven violation: at rate
\ResObsOobVtwoFpRate{} in the constructed study, where the share is seeded
outside the framework's record, and at rate \ResRsMrBookkeepingFpRate{} on
Gitea, where the grant is made between the crawl and the request. The
predicate's two representations cannot describe either grant. The error is a property of the
\emph{reading}, not of the system.
(2)~\emph{Oracle (RQ2).} When a readable, sound authorization oracle fully
resolves the proposition under test, does the metamorphic step add incremental
adjudication information beyond the oracle and the request outcome? Not
observably, within the evaluated scenarios: the step carries no additional
decision-relevant information beyond what the oracle-plus-request pair already
determines---an empirical premise about the evaluated relations, formalised as a
conditional (Proposition~\ref{lem:oracle-suff})---and the
paired methods disagree on only \ResAbDisagreementN{} records, all favoring
\texttt{direct\_observe} (false proofs \ResAbDirectObserveFP{}
vs.~\ResAbVThreeFP{}, recall \ResAbDirectObserveRec{} vs.~\ResAbVThreeRec{}; a
descriptive, point-estimate comparison, not a formal hypothesis test, no margin
pre-specified, not domination).
(3)~\emph{Measurement (RQ3).} Can a pooled specificity over follow-up inputs
provide stable scenario-level resolution for a failure confined to a class? No:
unaffected inputs dilute it, and a class the crawl never reaches contributes no
negative inputs at all, so the figure is silent about that class
(Proposition~\ref{prop:resolution}). The published labels are manually
verified; the limit is the denominator and the input set, not the labels.

The native execution, the seeded target and the second implementation are support
evidence for these findings, not further headline findings
(Section~\ref{sec:support}).

The findings are not a claim that \mbox{MST-wi} is wrong: the reachability
stand-in is announced by its authors and the predecessor design attributes its
residual false alarms to crawler coverage~\cite{Mai2019Metamorphic}. Our
contribution is a \emph{security-testing characterization} of metamorphic
testing for object-level authorization in Web/API systems, not a new detector,
with three layers. Although the technique is general, every empirical claim is
scoped to reachability-gated authorization relations in Web/API systems---a
prevalent class of networked-application security flaws---and we state where the
transfer stops. (1)~\emph{mechanistic}---the false-proof failure is located
structurally, in \emph{when} the reachability model was read rather than which
objects are visible; (2)~\emph{fair-control ablation}---under a readable, sound
oracle, a same-observable direct adjudicator shows the metamorphic step adds no
measurable adjudication information beyond the oracle and the request outcome;
and (3)~\emph{metric characterization}---a specific failure mode is identified in
a published specificity figure, which is not conditioned on the
reachability--authorization correspondence its interpretation assumes and whose
input-level denominator cannot register a class-uniform failure. The
evidence---the three-phase instantiation with internal controls
(Section~\ref{sec:protocol}), the scenario/configuration study
(Section~\ref{sec:synthetic}), the Gitea study with an endogenous control and a
seeded-vulnerability positive control (Section~\ref{sec:realsystem}), and the
source reading, the native execution of the library's own relations and the
second implementation (Section~\ref{sec:support})---supports these layers rather
than standing as further headline claims.

\section{Background: the predicates in question}
\label{sec:background}

For authorization there is no intrinsic oracle: vulnerable and correct
executions both answer with success. Metamorphic testing replaces the absent
oracle with a relation over multiple executions---and the relation is only as
sound as the assumptions it needs, which for authorization are exactly the facts
the oracle was supposed to supply. \mbox{MST-wi}~\cite{Chaleshtari2022Metamorphic},
building on Mai et al.~\cite{Mai2019Metamorphic}, crawls a per-user model of
reachable URLs and contents, then checks MRs against that recorded model. In the
artifact we inspected, the predicates guarding an authorization MR read either
the recorded model directly (\texttt{cannotReachThroughGUI},
\texttt{userCanRetrieveContent}) or static subject properties; we call a
predicate \emph{GUI-derived} when its value is a function of the crawl.

\paragraph{Why the false proof is a design consequence, not a bug} Suppose a
subject can reach an object at crawl time, so the record stores it as reachable.
Afterwards an administrator grants the same subject access to a \emph{further}
object through a channel---an administrative endpoint or a share link---that a
crawl of the subject's own interface never traverses. At execution the guard
still reads the frozen record, finds the new object absent, and concludes the
subject cannot reach it; the relation then treats the successful response as a
violation. Every step is exactly as specified, so no guard is miswritten: the
predicate's two representations---``reachable as recorded'' and ``not reachable
as recorded''---cannot express ``reachable, but only after the record was
taken''. More crawling does not remove the false proof, because the grant need
not be visible through the subject's interface at all, and the reported
specificity never registers it, because the error is a property of the
adjudication's \emph{reading moment} rather than of any single generated input.
This is the mechanism the rest of the paper measures, and the reason we speak of
a boundary of the specification rather than a defect of an implementation.

Two census figures frame what follows. Over the \ResCatalogMrTotal{} catalog
MRs, \ResCatalogGuiDerived{} (\ResCatalogGuiDerivedPct{}\%) gate on a
GUI-derived or identity predicate, and only \ResCatalogUnion{}
(\ResCatalogUnionPct{}\%) are authorization-relevant by name or action---an
upper bound, reported as a caveat; a full-text search for predicates reading the
\emph{state of an authorization artifact} (ownership, permission record, share)
finds \ResCatalogAuthzStateHits{} occurrences. Within the examined vocabulary,
then, the authorization MRs have no authorization-artifact state observable to
decide on, and the reading moment enters the specification: if an entitlement
changes after the crawl, the MR evaluates reachability against a world that no
longer exists. This is a statement about what the relations \emph{mean}; our
experiments are built so that the crawl-time and live readings can be told apart.

\section{Protocol: a faithful three-phase instantiation}
\label{sec:protocol}

Two instruments carry distinct claims. The three-phase instantiation below
carries the quantitative results of Sections~\ref{sec:synthetic}
and~\ref{sec:realsystem}; a direct execution of the library's own
relations---source unmodified, with the build configuration and a replay harness
added---carries one qualitative claim, that the published code produces the
false proof we predict (Section~\ref{sec:native-run}). Neither instrument is a
benchmark: the standing scope and non-claims are collected in
Section~\ref{sec:threats}, and neither yields a claim of detector accuracy.

\paragraph{Three phases}
(A)~\emph{Crawl}: record, per subject and object, whether the subject can reach
the object through the interface, and freeze that record---the model the
predicates will consult. (B)~\emph{Change authorization out of band}: alter a
(subject, object) entitlement through a channel outside the crawled interface;
the frozen model is not updated, because the architecture has no mechanism to
update it. (C)~\emph{Execute the MR} against the current state while the
predicates consult the frozen model. Phase~B is what the original evaluation
does not exercise, and it is where the reachability stand-in can be wrong
(Fig.~\ref{fig:protocol}).

\begin{figure*}[t]
\centering
\includegraphics[width=0.72\textwidth]{fig1_protocol.pdf}
\caption{The three-phase instantiation. The predicates consult the reachability
model frozen in Phase~A. A Phase~B entitlement change made outside the crawled
interface leaves that model stale, so in Phase~C a legitimately granted access
has no representation but ``cannot reach''; the exception predicate is not
discharged and the alarm is promoted to a proven violation---the false proof of
Finding~1.}
\label{fig:protocol}
\end{figure*}

\paragraph{Adjudicator variants and internal controls}
A raw alarm becomes a verdict through an exception policy, and different policies
are different methods: \textbf{v1} (raw alarms proven); \textbf{v2}
(legitimate-sharing exceptions filtered); \textbf{v3} (state-calibrated---consults
a readable authorization introspection surface, abstains when none is readable);
\texttt{degrade\_all} (the floor: abstain on every alarm);
\texttt{direct\_only} (authorization surface only, access outcome ignored---a
deliberately unfair control); and \texttt{direct\_observe} (a succeeding plain
request checked against the same readable surface). The last consumes the same
observable quantity as the MR family and decides the negative result.

\paragraph{Ground truth}
Every quantitative result names its ground-truth source, and it is never the
detector under test: construction state for the constructed system and the
scenario/configuration study; system-specific live re-probes of effective
authorization (the per-system validation procedure is described in
Section~\ref{sec:realsystem}) for Gitea, Gogs and GitLab; the pinned route tables for the grant census;
and the seed list for the seeded target.

\paragraph{Scenarios and design cells}
Seven scenarios separate the reasons an MR can misfire: no shared object (S1,
S2); an in-band legitimate share visible in the framework log (S3); an unshared
plain violation (S4); a read-only share plus a write violation (S5); an
out-of-band share invisible to the framework log (S6); an out-of-band share plus
a genuine violation (S7). They are crossed with \ResSynConfigs{} configurations
defined by instantiation mode, enforcement correctness, channel visibility and
introspection availability. The nominal cross is
$\ResDesignScenarios{}\times\ResSynConfigs{}$ combinations,
but four are inadmissible: the two no-exception configurations
(\texttt{V-noE3-vuln0}, \texttt{V-noE3-vuln1}) are undefined under the two
out-of-band scenarios (S6, S7), where the legitimate-sharing exception does not
apply. This leaves \ResDesignCells{} admissible design cells.

\paragraph{Statistical treatment}
The \ResSynRecords{} records are not independent observations. Each admissible
cell is expanded deterministically by the factors that are actually varied:
\ResDesignMrs{} metamorphic relations (MR-002, MR-004), \ResDesignPrecondVariants{}
exception-policy variants (raw, one level, two levels---the first deliberately
thrown switch), and \ResDesignGuiModes{} reachability readings (the frozen
crawl-time snapshot and the live execution-time state---the second deliberately
thrown switch). The two relations do not both apply to every cell: only MR-002 applies
under scenario~S5, so \ResDesignCellsOneMr{} S5 cells carry one relation while
the remaining \ResDesignCellsTwoMr{} carry two. The record count is therefore
\[
\ResSynRecords{} = \ResDesignCellsTwoMr{}\cdot\ResDesignMrs{}\cdot\ResDesignPrecondVariants{}\cdot\ResDesignGuiModes{}
+ \ResDesignCellsOneMr{}\cdot\ResDesignPrecondVariants{}\cdot\ResDesignGuiModes{},
\]
a deliberately unbalanced design---the per-cell multiplier differs by
scenario. The records are therefore a fixed enumeration, not a sample: we report
counts, rates and design-cell denominators, and we attach no inferential
interval to any of them. Where a headline count has a small denominator we say
so explicitly (the anchored rate, \ResFpEthreeHits{}/\ResFpEthreeTotal{}, is
unstable by construction) and lead with the larger-denominator count of the same
phenomenon. No conclusion rests on statistical significance. Where a comparison
is paired---the ablation of Section~\ref{sec:synthetic-ablation}---we report
discordant-pair counts rather than an interval over a difference. The
replication ceiling is
explicit: one constructed system, \ResDesignCells{} cells, \ResDesignMrs{} MRs,
three real systems.

\paragraph{Tooling and code provenance}
Development of the analysis scripts was assisted by an AI coding assistant
(declared in the Declarations). Every script that assistant touched was
reviewed, executed, and its output checked by the authors against the raw
records; the released pipeline recomputes every reported scalar from those
records, every table is generated by the pipeline, and the macro files and the
two auto-generated tables share a single source of truth. No AI tool
participated in the detection loop or in any adjudication.

\section{Synthetic study: where the adjudication goes}
\label{sec:synthetic}

This study supplies the evidence for Findings~1 and~3, and the ablation behind
Finding~2; the findings themselves are stated in Section~\ref{sec:findings}.

\subsection{Definitions}
\label{sec:metrics}

Confusions are defined over the \emph{verdict}: a \emph{false proof} is the
strict subset of false alarms promoted to \texttt{violation-proven}---an error
the tool asserts as established. We also report the \emph{detection-keep rate}
(share of fired true violations still alarmed) and count only abstentions
justified by an unreadable oracle.

\subsection{What the constructed system shows}

\paragraph{Attribution is clean} On the anchored configuration the false-proof
rate is \ResFpEthreeHits{}/\ResFpEthreeTotal{}. With that denominator the rate is
unstable by construction, so we attach to it no interval---nor a percentage,
which would travel out of
context. With the legitimate-sharing
exception absent the rate is \ResBaseFpRate{}
(\ResBaseFpHits{}/\ResBaseFpTotal{}), and with the channel listed it is
\ResListedCtlFpRate{} (\ResListedCtlFpHits{}/\ResListedCtlFpTotal{}). The
phenomenon is therefore attributable to a legitimate sharing exception the MR
cannot know about, not to execution itself. A $2{\times}2$ grid over channel
visibility and enforcement correctness tells the same story from the other
side: on unlisted cells the false-proof rate is
\ResBEThreevulnZerounlistedFpPct{}\% with no false negatives; on listed cells it
drops to \ResBEThreevulnZerolistedFpPct{}\% but one of the five violating records
in that cell is a false negative (\ResBEThreevulnOnelistedTp{} detected,
\ResBEThreevulnOnelistedFn{} missed). Avoiding false proofs and keeping detection are
controlled by one switch, and the switch cannot be thrown both ways at once.

\paragraph{Filtering exceptions does not help}
The exception ladder---raw, one exception level, two---leaves the anchored rate
at \ResLadderrawRate{}, \ResLaddereOneRate{}, \ResLaddereOneeTwoRate{}: it does
not move. The exception levels we can express are themselves functions of the
same reachability model, so they inherit its blind spot.

\paragraph{The framework cannot see the channel that matters}
When a legitimate share was created through the framework's facility, its log
holds \ResObsFwLogRecords{} records and v2 produces \ResObsFwVtwoFalseProof{}
false proofs (rate \ResObsFwVtwoFpRate{}); when the same share was created out
of band, the log holds \ResObsOobLogRecords{} records---no evidence the share
exists---and v2's false proofs rise to \ResObsOobVtwoFalseProof{} of
\ResObsOobFalseAlarm{} alarms (rate \ResObsOobVtwoFpRate{}), every one promoted
to \texttt{violation-proven}. The MR cannot distinguish ``this share does not
exist'' from ``this share exists but was not crawled''; its only representation
of the first is the second. On S7, where a genuine violation \emph{does} sit
behind the out-of-band share, the same blindness makes the alarm right for the
wrong reason (\ResObsOobBolaVtwoProvenBola{} of \ResObsOobBolaBolaTotal{}
reached)---and that coincidence of alarms with violations is what an
input-level metric cannot unpack. The same tension appears within a single
MR: for S3 (must not alarm) vs S5 (must alarm), no configuration in the family
is correct on both---the choice of switch is a choice of \emph{which} error to
make (Fig.~\ref{fig:scenarios}).

\begin{figure*}[t]
\centering
\includegraphics[width=0.78\textwidth]{fig2_scenarios.pdf}
\caption{The errors are scenario-structured. The bookkeeping oracle's false proofs
concentrate on~S6, the legitimate share created out of band and so invisible to
the framework's own log, while the state-calibrated v3 abstains in two places:
where its surface is not readable, and where a readable surface shows a covering
grant it declines to override. Both panels are recomputed from the same records
as Table~\ref{tab:ablation}.}
\label{fig:scenarios}
\end{figure*}

\paragraph{Stratification, not aggregation}
Aggregated, v3 looks best (recovery \ResVdVThreeRecoveryRate{} against
\texttt{degrade\_all}'s keep rate \ResVdDegradeallKeepRate{}); the reading is
wrong, because v3 changes regime with oracle readability. Stratified over the
\ResIntroOnN{}-record readable grid and its \ResIntroOffN{}-record absent grid (a
representative grid; Section~\ref{sec:stratum} re-runs the same comparison over
all \ResSynRecords{} records): on the readable grid it removes v2's false proofs
(\ResIntroOnVthreeFpN{} vs.~\ResIntroOnVtwoFpN{}) while still alarming on all
\ResIntroOnVthreeProvenBola{} fired true violations; on the absent grid
(\ResIntroOffN{}) it alarms on none of the \ResIntroOffVtwoProvenBola{} true
violations v2 still alarms on, abstaining \ResIntroOffAbstain{} times. Pooling
converts honest abstention into apparent performance, so v3's numbers may appear
only stratified.

\subsection{Fair ablation: does the metamorphic step add adjudication information?}
\label{sec:synthetic-ablation}

Our earlier iteration proposed v3 as a mechanism, resting on the claim that the
metamorphic step was necessary. We tested that claim by removing the step and
substituting the most naive consumer of the same observable. The distinction
matters: \texttt{direct\_only}---oracle yes, access outcome
ignored---produces \ResAbDirectOnlyFP{} false proofs against v3's
\ResAbVThreeFP{} (precision \ResAbDirectOnlyPrec{}), and rejecting it proves
nothing because it discards half the evidence. The fair control is
\texttt{direct\_observe}: a succeeding ordinary request plus the same readable
surface, consuming exactly what the MR-based adjudicator consumes.

\begin{table*}[t]
\centering
\caption{End-to-end confusion matrix over all \ResSynRecords{} records
(\ResAbNBola{} true violations, \ResAbNNormal{} correct cases). The records are
a fixed enumeration, not a sample (Section~\ref{sec:protocol}, statistical
treatment), so no inferential
interval is attached; where a denominator is small we say so rather than
interval it. \texttt{direct\_only} is a deliberately unfair control, shown to
make that point.}
\label{tab:ablation}
\begin{tabular}{lrrrrrrl}
\toprule
method & TP & FP & FN & TN & precision & recall & F1 \\
\midrule
raw (no exception handling) & \ResAbRawTP{} & \ResAbRawFP{} & \ResAbRawFN{} & \ResAbRawTN{} & \ResAbRawPrec{} & \ResAbRawRec{} & \ResAbRawFone{} \\
v1 (existence filter) & \ResAbVOneTP{} & \ResAbVOneFP{} & \ResAbVOneFN{} & \ResAbVOneTN{} & \ResAbVOnePrec{} & \ResAbVOneRec{} & \ResAbVOneFone{} \\
v2 (existence + coverage) & \ResAbVTwoTP{} & \ResAbVTwoFP{} & \ResAbVTwoFN{} & \ResAbVTwoTN{} & \ResAbVTwoPrec{} & \ResAbVTwoRec{} & \ResAbVTwoFone{} \\
v3 (state-calibrated) & \ResAbVThreeTP{} & \ResAbVThreeFP{} & \ResAbVThreeFN{} & \ResAbVThreeTN{} & \ResAbVThreePrec{} & \ResAbVThreeRec{} & \ResAbVThreeFone{} \\
\texttt{degrade\_all} & \ResAbDegradeAllTP{} & \ResAbDegradeAllFP{} & \ResAbDegradeAllFN{} & \ResAbDegradeAllTN{} & \ResAbDegradeAllPrec{} & \ResAbDegradeAllRec{} & \ResAbDegradeAllFone{} \\
\texttt{direct\_only} & \ResAbDirectOnlyTP{} & \ResAbDirectOnlyFP{} & \ResAbDirectOnlyFN{} & \ResAbDirectOnlyTN{} & \ResAbDirectOnlyPrec{} & \ResAbDirectOnlyRec{} & \ResAbDirectOnlyFone{} \\
\texttt{direct\_observe} & \ResAbDirectObserveTP{} & \ResAbDirectObserveFP{} & \ResAbDirectObserveFN{} & \ResAbDirectObserveTN{} & \ResAbDirectObservePrec{} & \ResAbDirectObserveRec{} & \ResAbDirectObserveFone{} \\
\bottomrule
\end{tabular}
\end{table*}

Table~\ref{tab:ablation} concludes the ablation. The first three rows are a
ladder, and they must be read in order: with no exception handling at all,
\texttt{raw} yields \ResAbRawFP{} false proofs; the existence filter v1 cuts
them to \ResAbVOneFP{}, but pays for it by abstaining on true violations (FN
rises to \ResAbVOneFN{}); adding operation coverage (v2) restores the true
violations (FN \ResAbVTwoFN{}) at the same false-proof count. Neither filter is
therefore a free improvement, and v3's zero in the FP column is bought with the
FN column. \texttt{direct\_observe} has
\ResAbDirectObserveFP{} false proofs against v3's \ResAbVThreeFP{}---a tie at
zero, not a win. The comparison is paired, so we report the discordance rather
than an inferential interval: over the \ResAbNBola{} true-violation records the
two methods disagree on \ResAbDisagreementN{}, and on every one of them
\texttt{direct\_observe} is right---it detects a violation that v3 misses or
abstains on, its reachability precondition blocking a detection the plain
request makes. The aggregate recalls \ResAbDirectObserveRec{} and
\ResAbVThreeRec{} restate that asymmetry rather than measure a difference, and
no margin was pre-specified, so no formal hypothesis-test claim attaches. We
found no incremental adjudication advantage of the metamorphic step over direct
observation in these oracle-available scenarios, and that absence provides no
support for an incremental adjudication benefit in these scenarios:
the step would have to be better at something the plain request cannot do.
Selective presentation would have hidden the cost: v3's precision of
\ResAbVThreePrec{} was bought with \ResAbVThreeFN{} false negatives (F1
\ResAbVThreeFone{}), an exchange of \ResAbVOneFP{} false alarms for missed
violations dressed as an improvement.

\subsection{Coverage, selective risk, and what a pooled precision does not say}
\label{sec:coverage}

Table~\ref{tab:ablation} scores the alarm. A method that abstains is counted as
not having alarmed, so the abstention lands in TN or in FN: for the question
``does the tool report a violation'' that is the right scoring, but it is not the
scoring for ``does the method decide''. Table~\ref{tab:coverage} rescores the
same \CvgN{} records on the decision and reports coverage, abstention rate,
selective risk and request cost.

v3 decides \CvgVthreeDecided{} of \CvgN{} records---coverage
\CvgVthreeCoveragePct{}\%---and abstains on \CvgVthreeAbstain{},
\CvgVthreeAbstBola{} of them on true violations and \CvgVthreeAbstNormal{} on
records with none. Its precision of $1.00$ and its selective risk of
\CvgVthreeRiskPct{}\% are therefore properties of a \CvgVthreeCoveragePct{}\% coverage: over the
remaining two records in five the method claims nothing, and the pooled precision
is not an error rate over all records.

The abstentions are not all of one kind (Table~\ref{tab:coverage}, panel b).
\CvgVthreeUnread{} are structural---the surface returns nothing, and there is
nothing for the method to consult. The other \CvgVthreeCover{} fall on a
\emph{readable} surface that shows a covering grant: the access is explained, and
the method still declines to assert its absence. A conservative method may prefer
silence to a positive ``no violation'', and we do not score that as an error; but
it is a non-decision, and it cannot be read as a correct rejection. The
comparison with v2 makes the arithmetic visible. v2 promotes \CvgVtwoFP{} false
proofs; v3 declines \CvgVtwoFPUnread{} of them because the surface is unreadable
and \CvgVtwoFPCover{} because it is readable and shows the grant. Only the first
\CvgVtwoFPUnread{} are the availability gate the method is named for. The rest is
a false proof converted into a non-decision.

The other side of the ledger is dearer. \CvgVthreeAbstBola{} abstentions fall on
true violations, and with the \CvgVthreeMissed{} records on which no alarm is
raised at all, v3 leaves \CvgVthreeFNunproven{} of \CvgViolations{} violations
unproven---the count the review derived from the pooled table.

\texttt{direct\_only} is scored under its own adjudication rule, which is also
the rule Table~\ref{tab:ablation} uses: assert a violation whenever the
authorization surface reports no granting channel, regardless of the request
outcome. On the \CvgN{} records it decides \CvgDirOnlyDecided{} (the same
abstention set as \texttt{direct\_observe}, so the two differ in rule, not in
coverage) and it is right on \CvgDirOnlyTP{} true violations. Its
\CvgDirOnlyFP{} false proofs are not, however, strict false proofs in the sense
defined in Section~\ref{sec:protocol}: every one of them falls on a request the
server in fact denied, so no unauthorized access resulted
(\CvgDirOnlyFpNoEffect{} of \CvgDirOnlyFP{} without violation effect,
\CvgDirOnlyFpEffect{} on allowed requests). They are invalid violation
claims---the oracle's ``no granting channel'' read as a violation claim that the
request outcome contradicts---and its selective risk over the decided set is
\CvgDirOnlyRiskPct\%, not zero. Keeping the rule fixed across both tables is what
exposes this cost; a table that silently added the request-success condition to
this row would have reported a zero risk that the method never earned.
The overlap of v3's decisions with \texttt{direct\_only}'s covers
\CvgJointVthreeDirect{} records---the \CvgObsDecided{} records on which the direct
method decides minus the \CvgVthreeCover{} records on which v3 additionally
declines on a readable covering surface---and the violation-detection difference
on comparable denominators is stated with its own denominator, not with the
\CvgObsDecided{} oracle-readable count.

\texttt{direct\_observe} is scored on the same contract, and that changes what
its earlier numbers meant. It abstains when the probe is unreadable, exactly as
the direct adjudicator does, so it decides \CvgObsDecided{} records---coverage
\CvgObsCoveragePct{}\%---and abstains on \CvgObsAbstain{}, \CvgObsAbstBola{} of
them true violations. Within the records it decides it detects
\CvgObsDetected{} and raises no false proof. On the \CvgJointVthreeDirect{} records
both methods decide, it detects \CvgDetectGap{} more violations than v3; the
remaining difference in pooled recall is coverage, not accuracy. The last row of
Table~\ref{tab:coverage} is the binary version that was implicit before: it
forces a negative verdict where the probe is unreadable, covers all
\CvgN{} records, detects the same \CvgObsDetected{}, and misses
\CvgForcedObsFN{}. Its selective risk of \CvgForcedObsRiskPct{}\% is entirely
missed violations, produced by the forcing rather than by the method; we report
it as a reference point, not as a comparator.

Selective risk is not a ranking either, and we do not use it as one.
\texttt{degrade\_all}, which alarms on nothing, has a selective risk of
\CvgDegradeRiskPct{}\% at \CvgDegradeCoveragePct{}\% coverage and detects no
violation at all; an ordering by risk alone would place it above v2. We therefore
report precision and selective risk with coverage attached and never compare them
across methods at unequal coverage. The stratified analysis
(Section~\ref{sec:stratum}) equalises the \emph{availability} of the oracle, not
the decision rate: v3 still declines \CvgVthreeCover{} readable records that show
a covering grant, so the two strata do not make coverage common, and we say so
wherever the comparison is made. A safe abstention is not a detection and we
never count it as one.

\paragraph{Cost} The enumeration issues \CvgReqTotal{} requests over \CvgN{}
records (median \CvgReqMedian{} per record, range \CvgReqMin--\CvgReqMax). The
oracle-consuming methods add \CvgProbeGets{} probe GETs over \CvgProbeCells{}
cell-operations: charged per record that is one request, and \CvgProbePerRecord{}
amortised as executed, because the probe is shared by the records of a cell.
\texttt{raw}, v1 and v2 issue none of their own---they read the log the crawl
already wrote. We do not report wall-clock time for this study: every method
replays the same recorded exchanges, so a difference would measure our replay
loop rather than the system under test. The study we did execute end to end is
the temporal loop of Section~\ref{sec:realsystem}: a median of
\CvgLoopGiteaTotal{}~s per transition on Gitea and \CvgLoopGitlabTotal{}~s on
GitLab, of which the mutation is \CvgLoopGiteaMut{}~s and \CvgLoopGitlabMut{}~s,
the recrawl \CvgLoopGiteaRecrawl{}~s and \CvgLoopGitlabRecrawl{}~s, and the
cross-credential probe sweep \CvgLoopGiteaProbe{}~s and \CvgLoopGitlabProbe{}~s.
Each of these is a separate median, so the parts need not sum to the total
exactly. The recrawl the mechanism needs is a small part of the loop; re-reading
every credential is what costs.

\input{coverage_table.tex}

\subsection{A published metric's resolution: what a pooled specificity cannot register}
\label{sec:resolution}

The published figure is a specificity, not a false-alarm share of all inputs.
\mbox{MST-wi} defines it as $\mathrm{TN}/(\mathrm{TN}+\mathrm{FP})$ over the
follow-up inputs $F$, and states that each reported MT failure was manually
verified~\cite{Chaleshtari2022Metamorphic}; the labels are therefore not the
crawler's own answer. With $\mathrm{FN}=0$ throughout that evaluation,
$\mathrm{TN}=F-\mathrm{TP}-\mathrm{FP}$ and the figure equals
$1-\mathrm{FP}/(F-\mathrm{TP})$: the denominator is the \emph{negative}
follow-up inputs, not every generated input. We use that form below, and we do
not attribute the published number to a circular label.

What such a figure cannot do is register a failure confined to a class, and the
two reasons are arithmetic and coverage rather than labelling. First, the figure
is pooled: a class contributing a share $w$ of the negative inputs can move it
by at most $w$, so a real class-uniform failure is diluted by unaffected
inputs. Second, the follow-up inputs are those the crawl could enumerate; a
class of legitimate access the crawl never reaches contributes no negative
inputs and therefore cannot enter the ratio at all. The second point is the
measurement face of Finding~1, and it is a statement about the coverage of the
input set, not about the correctness of the human labels. We cannot reproduce
the engine's exact input count, so we substitute no value for the published
figure; the argument below is about the form of the denominator only.

\paragraph{A bounded formal remark on the denominator}
Write $S=1-\mathit{FP}/N^{-}$, where $N^{-}=\mathrm{TN}+\mathrm{FP}$ is the
number of negative follow-up inputs the run generated, and define
$S_{\mathrm{scenario}} = 1 - N_{\mathrm{failed}} / N_{\mathrm{scenarios}}$,
where $N_{\mathrm{failed}}$ is the number of scenarios in which the failure is
uniform. $S_{\mathrm{scenario}}$ is an analytical construct for characterising
resolution; it is not proposed as a replacement metric. The remark concerns
resolution, not estimator consistency or power.

\begin{proposition}
\label{prop:resolution}
Fix $f \ge 1$ failed scenarios, each producing a fixed number $a \ge 1$ of
false-alarm inputs, and $s \ge 1$ successful scenarios each contributing $m$
additional alarm-free negative inputs. Then
$S = 1 - (f a) / (f a + s m)$, which tends to $1$ as $m \to \infty$, while
$S_{\mathrm{scenario}} = 1 - f/(f+s)$ is independent of $m$ and strictly less
than $1$.
\end{proposition}

\begin{proof}
$N^{-} = f a + s m$ and $\mathit{FP} = f a$; substitution gives the first
expression, which increases monotonically in $m$ to $1$, whereas the
scenario-level figure depends only on the fixed counts $f$ and $s$.
\end{proof}

\begin{corollary}
\label{cor:resolution}
No fixed input-level specificity threshold $t < 1$ can guarantee that a
scenario-uniform failure will push the metric below it: choosing
$m > (f a/s)\cdot t/(1-t)$ drives $S$ above $t$ while the $f$ failing
scenarios remain failing.
\end{corollary}

The finding is therefore about what the published figure can be asked to show,
not about its arithmetic: it is pooled over classes, and it is computed on the
inputs the crawl reached. Neither point disputes the reported value or its
labels; both limit what a high value certifies about a class the run did not
exercise.

\paragraph{Implementation robustness}
An independently written implementation of the same predicates was compared on a
subset fixed in advance; it bounds implementation accident only, and is reported
with the confirmations in Section~\ref{sec:secondimpl}.

% =====================================================================
%  M2 补充分析：把池化消融表按「oracle 是否有效可读」分层
%  重要：本节所有数字一律引用 m2_stratum.tex 的 \Mii* 宏，
%       由 analysis/ingest_m2_stratum.py 从 raw_matrix.jsonl 重算得到。
%       禁止手敲数字。
% =====================================================================

\subsection{What pooling hid: the ablation split by oracle readability}
\label{sec:stratum}

Table~\ref{tab:ablation} pools two populations in which the oracle-consuming
methods behave in opposite ways, and the pool is misleading in both directions.
Section~\ref{sec:synthetic} reports the comparison on a representative grid; this
subsection, summarised in Table~\ref{tab:stratum}, repeats it over all
\MiiNTotal{} records. The records differ in one respect that decides whether
those methods can act at
all: whether the introspection surface returns a reading for the record. We
therefore repeat the same comparison inside each stratum, with the same methods
and the same decision rule, and we add two references that need no oracle
against which the metamorphic family can be scored where no oracle exists.

\paragraph{Stratification variable} We stratify on the \emph{effective}
readability of the surface---the value the probe returns for the record---not on
the configuration switch. The distinction is not cosmetic: \MiiFlagOnButUnavailable{}
records sit in configurations flagged \texttt{introspection=on} whose surface
nevertheless returns nothing (the two no-exception configurations), so the
configuration flag is not a faithful proxy for whether an oracle was available
to the adjudicator, and using it would misclassify those records. The two strata
partition all \MiiNTotal{} records: \MiiNOn{} readable, \MiiNOff{} unreadable.
Both are fixed enumerations, so as elsewhere we report counts and rates and
attach no inferential interval.

\begin{table*}[t]
\centering
\caption{The ablation of Table~\ref{tab:ablation} re-run inside each stratum.
The two panels are not rankings of the same population but two different
regimes: in the readable stratum the oracle-bearing methods can act; in the
unreadable stratum they cannot act and are structurally bound to zero
detections. The two reference rows need no oracle and no crawl model:
\texttt{always\_alarm} is the trivial floor, \texttt{obs\_success} alarms on the
observed success of a plain request. All numbers are counts.}
\label{tab:stratum}
\begin{tabular}{llrrrrrrl}
\toprule
stratum & method & TP & FP & FN & TN & precision & recall & F1 \\
\midrule
readable (\MiiNOn)
  & \texttt{raw}               & \MiiOnRawTP & \MiiOnRawFP & \MiiOnRawFN & \MiiOnRawTN & \MiiOnRawPrec & \MiiOnRawRec & \MiiOnRawFone \\
  & v1 (existence filter)   & \MiiOnVOneTP & \MiiOnVOneFP & \MiiOnVOneFN & \MiiOnVOneTN & \MiiOnVOnePrec & \MiiOnVOneRec & \MiiOnVOneFone \\
  & v2 (existence+coverage) & \MiiOnVTwoTP & \MiiOnVTwoFP & \MiiOnVTwoFN & \MiiOnVTwoTN & \MiiOnVTwoPrec & \MiiOnVTwoRec & \MiiOnVTwoFone \\
  & v3 (state-calibrated)   & \MiiOnVThreeTP & \MiiOnVThreeFP & \MiiOnVThreeFN & \MiiOnVThreeTN & \MiiOnVThreePrec & \MiiOnVThreeRec & \MiiOnVThreeFone \\
  & \texttt{direct\_observe}& \MiiOnDirectObserveTP & \MiiOnDirectObserveFP & \MiiOnDirectObserveFN & \MiiOnDirectObserveTN & \MiiOnDirectObservePrec & \MiiOnDirectObserveRec & \MiiOnDirectObserveFone \\
  & \texttt{obs\_success}   & \MiiOnObsSuccessTP & \MiiOnObsSuccessFP & \MiiOnObsSuccessFN & \MiiOnObsSuccessTN & \MiiOnObsSuccessPrec & \MiiOnObsSuccessRec & \MiiOnObsSuccessFone \\
  & \texttt{always\_alarm}  & \MiiOnAlwaysAlarmTP & \MiiOnAlwaysAlarmFP & \MiiOnAlwaysAlarmFN & \MiiOnAlwaysAlarmTN & \MiiOnAlwaysAlarmPrec & \MiiOnAlwaysAlarmRec & \MiiOnAlwaysAlarmFone \\
\midrule
unreadable (\MiiNOff)
  & \texttt{raw}               & \MiiOffRawTP & \MiiOffRawFP & \MiiOffRawFN & \MiiOffRawTN & \MiiOffRawPrec & \MiiOffRawRec & \MiiOffRawFone \\
  & v1 (existence filter)   & \MiiOffVOneTP & \MiiOffVOneFP & \MiiOffVOneFN & \MiiOffVOneTN & \MiiOffVOnePrec & \MiiOffVOneRec & \MiiOffVOneFone \\
  & v2 (existence+coverage) & \MiiOffVTwoTP & \MiiOffVTwoFP & \MiiOffVTwoFN & \MiiOffVTwoTN & \MiiOffVTwoPrec & \MiiOffVTwoRec & \MiiOffVTwoFone \\
  & v3 (state-calibrated)   & \MiiOffVThreeTP & \MiiOffVThreeFP & \MiiOffVThreeFN & \MiiOffVThreeTN & \MiiOffVThreePrec & \MiiOffVThreeRec & \MiiOffVThreeFone \\
  & \texttt{direct\_observe}& \MiiOffDirectObserveTP & \MiiOffDirectObserveFP & \MiiOffDirectObserveFN & \MiiOffDirectObserveTN & \MiiOffDirectObservePrec & \MiiOffDirectObserveRec & \MiiOffDirectObserveFone \\
  & \texttt{obs\_success}   & \MiiOffObsSuccessTP & \MiiOffObsSuccessFP & \MiiOffObsSuccessFN & \MiiOffObsSuccessTN & \MiiOffObsSuccessPrec & \MiiOffObsSuccessRec & \MiiOffObsSuccessFone \\
  & \texttt{always\_alarm}  & \MiiOffAlwaysAlarmTP & \MiiOffAlwaysAlarmFP & \MiiOffAlwaysAlarmFN & \MiiOffAlwaysAlarmTN & \MiiOffAlwaysAlarmPrec & \MiiOffAlwaysAlarmRec & \MiiOffAlwaysAlarmFone \\
\bottomrule
\end{tabular}
\end{table*}

\paragraph{The readable stratum strengthens Finding~2}
Splitting does not rescue the metamorphic step; it removes the comparison in
which it looked best. On the \MiiNOn{} readable records
\texttt{direct\_observe} reaches recall \MiiOnDirectObserveRec{} with
\MiiOnDirectObserveFP{} false proofs, against v2's recall \MiiOnVTwoRec{} with
\MiiOnVTwoFP{} false proofs, and against v3's recall \MiiOnVThreeRec{} with
\MiiOnVThreeFP{}. The ordering is therefore not an artifact of the metric---the
plain request is ahead on false proofs and not behind on detections---whereas
the pooled table had let the family appear competitive on recall. We do not
claim domination: the comparison rests on the constructed-observable argument of
Proposition~\ref{lem:oracle-suff}, and \texttt{direct\_observe} is not an independent
method here, since it is built from the same two quantities that define the
reference (its second input, the access outcome, is part of how the reference
labels a violation). What the split shows is the scope of the necessity claim:
where an oracle resolves the proposition, there is nothing left for the relation
to add, and the pooled numbers were flattering to the relation only because they
mixed in records where the oracle-bearing methods could not run.

\paragraph{The unreadable stratum is the relation's own regime, and it is not empty}
The \MiiNOff{} unreadable records contain \MiiOffNBola{} true violations and
\MiiOffNNormal{} correct cases, a violation base rate of
\MiiOffPrevalencePct{}\%. Every oracle-bearing method---v3,
\texttt{direct\_observe}, \texttt{direct\_only}---detects \MiiOffVThreeTP{} of
them, not because the relations are weak but because those methods have nothing
to consult and, by their own definition, abstain. The metamorphic family does
not abstain: v2 reaches all \MiiOffVTwoTP{} true violations
(recall \MiiOffVTwoRec{}) with \MiiOffVTwoFP{} false proofs, precision
\MiiOffVTwoPrec{}, a lift of \MiiLiftVTwo{}$\times$ over the base rate. Against
the two no-oracle references: \texttt{always\_alarm} scores precision
\MiiOffAlwaysAlarmPrec{}, which is the base rate itself---no information at
all---and \texttt{obs\_success}, which alarms on the observed success of a plain
request, reaches the \emph{same} recall at precision \MiiOffObsSuccessPrec{},
lift \MiiLiftObsSuccess{}$\times$. The net difference between the metamorphic
family and that reference is therefore not in detections
(\MiiDeltaTPVTwoMinusObsSuccess{} additional true violations) but in false
alarms: \MiiOffAttribSavedFPTotal{} fewer for the metamorphic family.

\paragraph{What the saving consists of}
The attribution is one-sided and mechanical. Over the seven scenarios, every
record of the \MiiOffAttribSavedFPTotal{}-alarm saving falls on S3, the
legitimate share created \emph{through the framework's own facility} and hence
recorded in its log: \MiiOffScenThreeObsSuccessFP{} false alarms for
\texttt{obs\_success} against \MiiOffScenThreeVTwoFP{} for v2 on the
\MiiOffScenThreeN{} records of that scenario. On S6, the same share created out
of band, the two agree exactly (\MiiOffScenSixObsSuccessFP{} vs.\
\MiiOffScenSixVTwoFP{} false alarms on \MiiOffScenSixN{} records), and both for the reason developed in
Section~\ref{sec:boundary1}---the predicate has no representation for a grant the
crawl did not see. So the incremental detection
information the metamorphic step carries where no oracle exists is exactly the
information the crawl-derived reachability model can see: it suppresses what was
logged, and it reproduces Finding~1 on what was not. This bounds Finding~2
rather than overturning it. Finding~2's antecedent---a readable, sound oracle
that fully resolves the proposition---fails on this stratum by construction, and
where it fails the relation is not worthless but worth precisely the reachability
model it reads: above a trivial floor, above a single uninterpreted request, and
below any method that can consult the entitlement.

\paragraph{Limits of the comparison}
Two references are deliberately weak, and we do not present them as the frontier
of oracle-free detection: \texttt{obs\_success} uses one request and no model,
and a differential prober that compares several accounts against one another---the
shape taken by the black-box line reviewed in
Section~\ref{sec:related}---consumes more than \texttt{obs\_success} does and was
not instantiated here. The lift we measure is therefore an estimate against weak
references, not a claim about metamorphic testing versus the best oracle-free
detector, and we make no such claim. Both references also share an observable
with the reference labelling (\texttt{obs\_success} alarms on the same access
outcome that partly defines a violation), so their absolute rates carry the same
constructiveness caveat as \texttt{direct\_observe}; only the differences are
meant to be read. Everything here is one constructed system, one MR family, and
one criterion (the alarm) rather than the promoted verdict.

\section{Real systems: Gitea, Gogs, GitLab, and a seeded control}
\label{sec:realsystem}

This section supplies the real-system evidence for Findings~1 and~2: the Gitea
study with its endogenous control, the cross-vendor replication, and the seeded
positive control (support evidence, Section~\ref{sec:support}).

A synthetic system cannot tell us whether the crawl-time reading is an artifact
of our construction, so we repeated the essential comparison on real systems. We
chose Gitea \ResRsVersion{} for genuine object-level sharing with per-object
visibility, a machine-readable sharing surface, and a documented administrative
grant path. Three scope conditions, stated before the numbers. First, on the
endpoints and authorization transitions exercised here, Gitea enforced the
expected decisions: no instance of the studied vulnerability was found, so the
experiment measures only the false-proof side. Second, the cell study uses our
instantiation of the predicates, not the engine; the engine's own relations are
executed separately in Section~\ref{sec:native-run}. Third, the reachability
model is reported under \emph{both} readings---crawl-time snapshot (what the
definitions presuppose) and the live execution-time state---so the time at which
the model was read can be isolated as the explanatory variable. The experiment
produces \ResRsCells{} cells, \ResRsCellsSnapshot{} per reading.

\begin{figure*}[t]
\centering
\includegraphics[width=0.86\textwidth]{fig5_real_system.pdf}
\caption{Gitea \ResRsVersion{}. (a)~the grant-API space; (b)~authorization-surface
readability, complete from the owner/admin vantage
(\ResRsOwnerReadable{}/\ResRsOwnerReadableTotal{}) and never from a black-box
subject's (\ResRsBlackboxSelfReadable{}/\ResRsBlackboxSelfCells{}); (c)~the
departure from the real authorization state is present under the crawl-time
reading (\ResRsMismatchSnapshot{}/\ResRsMismatchDenom{}) and absent live; (d)
alarms versus false proofs per adjudicator. All counts are recomputed from the
run records by the analysis pipeline.}
\label{fig:real-system}
\end{figure*}

\paragraph{Grant surface}
Table~\ref{tab:rs-m1} counts an endpoint as grant-capable if, without modifying
the system, a request can move a (subject, object) pair from not-entitled to
entitled. The \ResRsGrantNoAdmin{} endpoints drivable by a non-administrator are
the out-of-band channel in its real form: a subject who is not an administrator
can create an entitlement that a crawl of that subject's own interface need not
reveal. The cross-vendor table below reports a normalized census (semantically
equivalent routes merged, non-authorization grant surface excluded), which is why
its Gitea figure (\GtGrantEndpoints{}) differs from the raw count
(\ResRsGrantEndpoints{}) here; the two answer different questions.

\begin{table}[t]
\centering
\caption{The grant surface of Gitea \ResRsVersion{} (instance-independent).}
\label{tab:rs-m1}
\begin{tabular}{lr}
\toprule
quantity & count \\
\midrule
paths in the system's API description & \ResRsSwaggerPaths{} \\
grant endpoints & \ResRsGrantEndpoints{} \\
\quad of which usable without site-admin & \ResRsGrantNoAdmin{} \\
revocation endpoints & \ResRsRevokeEndpoints{} \\
\bottomrule
\end{tabular}
\end{table}

\paragraph{Is the authorization surface readable, and from where?}
Readability depends entirely on the vantage, so the counts below are different
populations and must not be read as one denominator. From the owner/administrator
vantage the sharing surface is readable on \ResRsOwnerReadable{} of
\ResRsOwnerReadableTotal{} owner-vantage attempts over the
\ResRsCellsSnapshot{} cells: complete. From a black-box subject's own
vantage---the vantage a black-box tool occupies---it is readable on
\ResRsBlackboxSelfReadable{} of \ResRsBlackboxSelfCells{}: never. An intermediate
count of \ResRsSelfReadable{}/\ResRsSelfTotal{} comes only from subjects who
happen to be owners or administrators and does not constrain the mechanism; it is
not synonymous with ``an arbitrary black-box subject can read the surface''.
Table~\ref{tab:introspection-readable} collects these populations. On
\ResRsFidelityConsistent{} of \ResRsFidelityTotal{} applicable cells the
credential-substitution check confirmed that our model agrees with the system's
behavior (Fig.~\ref{fig:real-system}).

\paragraph{Is the write path verified by state rather than status?}
A status code cannot establish that a write took effect, so the write path is
checked by content: the subject creates a uniquely named file, the owner
re-reads it byte-for-byte as the ground truth, and the file is deleted where the
API allows it. Across the \ResWpSystems{} systems, \ResWpAgreeTotal{} of
\ResWpPairsTotal{} cross-credential write attempts agree with the ledger, with
\ResWpUnauthTotal{} unauthorized writes; no status code contradicted the content
outcome (count: \ResWpStatusOnlyDisagreements{}). The owner control---a non-admin
writing its own repository---changed content in \ResWpOwnerObservedTotal{} of
\ResWpOwnerRowsTotal{} attempts. On \ResWpAdminBlockedSystems{} of the
\ResWpSystems{} systems the site administrator could not write through the
contents API although the ledger grants write; those \ResWpAdminRowsTotal{}
control rows are excluded from the cross-credential denominator and reported as
a modelling gap rather than absorbed. Restore is unavailable on
\ResWpRestoreUnsupported{} of the \ResWpSystems{} systems, whose contents API
exposes no delete, so its probe files remain in a disposable instance.

\begin{table}[t]
\centering
\caption{Authorization-introspection readability protocol on Gitea. Each row is a
distinct population and denominator, so the rows are not comparable: both the
vantage and the unit of the denominator differ.}
\label{tab:introspection-readable}
\small
\begin{tabular}{@{}p{0.44\columnwidth}rr@{}}
\toprule
vantage & attempts & readable \\
\midrule
owner-vantage authorization-introspection surface & \ResRsOwnerReadableTotal{} & \ResRsOwnerReadable{} \\
arbitrary subject (black box) & \ResRsBlackboxSelfCells{} & \ResRsBlackboxSelfReadable{} \\
subject also owner/admin (representative grid) & \ResRsSelfTotal{} & \ResRsSelfReadable{} \\
\bottomrule
\end{tabular}
\end{table}

\paragraph{The mismatch is in time, not in visibility}
Table~\ref{tab:rs-m4} is the load-bearing result: the adjudication departs from
the real authorization state on \ResRsMismatchSnapshot{} of \ResRsMismatchDenom{}
cells under the crawl-time reading and \ResRsMismatchLive{} of the same cells
under the live reading. System state and tested relation were held fixed; only
the time at which the reachability model was read changed. Under the crawl-time
reading a bookkeeping adjudication produces \ResRsMrBookkeepingFp{} false proofs
on the \ResRsAuthorizedCells{} authorized cells (rate
\ResRsMrBookkeepingFpRate{}), against \ResRsMrBookkeepingAlarmsLive{} alarms
live; v3 produces \ResRsVThreeFp{} false proofs, and \texttt{direct\_observe}
emits \ResRsDirectObserveFp{}. The unfair control \texttt{direct\_only} raises
\ResRsDirectOnlyAlarms{} alarms, more than any other adjudicator, simply because
it ignores the request outcome; on this system none is promoted to
\texttt{violation-proven} under our strict false-proof definition
(\ResRsDirectOnlyFp{} false proofs), so alarm count and false-proof count are not
the same quantity. Channel provenance of the authorized access
(\ResRsChannelNotObservable{} of \ResRsChannelAuthorizedTotal{} not observable
by the framework log, \ResRsChannelNotObservableRate{}; of these
\ResRsChannelOobLiteral{} are literal out-of-band grants and
\ResRsChannelImplicit{} implicit via site administration or public visibility) is
a property of our
topology---a design parameter, not a field prevalence. This comparison uses two
snapshots and therefore cannot separate a stale model from a channel the crawl
never listed; the transition experiment in Section~\ref{sec:temporal} separates
the two.

\begin{table}[t]
\centering
\caption{Departure from the real authorization state, per reading of the
reachability model.}
\label{tab:rs-m4}
\begin{tabular}{@{}p{0.56\columnwidth}rr@{}}
\toprule
reading of the reachability model & mismatch & cells \\
\midrule
crawl-time snapshot (as the definitions presuppose) & \ResRsMismatchSnapshot{} & \ResRsMismatchDenom{} \\
live execution-time state & \ResRsMismatchLive{} & \ResRsMismatchDenom{} \\
\bottomrule
\end{tabular}
\end{table}

\paragraph{An endogenous control}
The most serious alternative explanation is that public objects are simply harder
to enumerate---visibility, not time. Two objects that end equally public but
became public at different moments control for this (Table~\ref{tab:rs-endo}):
the predicate reports unreachability on \ResRsEndoAfterCrawlReach{} of
\ResRsEndoAfterCrawlCells{} cells only for the object made public after the
crawl, and on \ResRsEndoFromStartReach{} for the object public throughout, with
the MR firing correspondingly (\ResRsEndoAfterCrawlFires{}
vs.~\ResRsEndoFromStartFires{}). Final visibility is identical, so within this
design the only manipulated explanatory variable is when the model was read.

\begin{table}[t]
\centering
\caption{Endogenous control: two objects that are both public at the end.}
\label{tab:rs-endo}
\begin{tabular}{@{}p{0.40\columnwidth}rrr@{}}
\toprule
object & cells & \texttt{cannotReach} & MR fires \\
\midrule
public before the crawl & \ResRsEndoFromStartCells{} & \ResRsEndoFromStartReach{} & \ResRsEndoFromStartFires{} \\
public only after the crawl & \ResRsEndoAfterCrawlCells{} & \ResRsEndoAfterCrawlReach{} & \ResRsEndoAfterCrawlFires{} \\
\bottomrule
\end{tabular}
\end{table}

\subsection{Eight committed transitions with a frozen model, on two products}
\label{sec:temporal}

The two-snapshot comparison above cannot separate two explanations: the model
went stale, or the crawl never listed the channel in the first place. The
transition experiment separates them by construction. We ran eight
authorisation transitions on one fresh Gitea \TmpVersion{} instance and repeated
the same eight, under the same pre-registered protocol (hash \texttt{\TmpProtocol})
and the same ledger rules, on an independent GitLab CE \TmpGVersion{} instance.
The mechanisms differ per product---attaching a repository to a team becomes
sharing a project with a group; removing a team member becomes removing a group
member---so the two runs share a protocol, not a topology, and we report them
side by side rather than pooled.

Each transition is a closed loop. We crawl every subject, serialise the
per-subject reachable set, freeze it as $M_0$ and record its digest; commit one
authorisation transaction and increment the harness version $hv$ only if the
transaction returns 2xx; re-crawl with the \emph{same} credentials to obtain
$M_1$; read the oracle from the owner vantage; issue the request and verify the
effect by content; and classify against the ledger at decision time. All
\TmpBothPolicyEvents{} transactions committed, all \TmpBothMzeroVerified{} frozen
$M_0$ digests were unchanged when re-checked immediately before probing, and all
\TmpBothVersionConsistent{} cycles were version-consistent. Across
\TmpBothCells{} cross-credential cells (\TmpCells{} on Gitea plus \TmpGCells{} on
GitLab), $E$ and $Y$ agree in \TmpBothEeqY{} cases, so no cell is scored on a
modelling discrepancy---under the implementation's closed-world default-deny
reading of $E$; Table~\ref{tab:unknown-sensitivity} reconciles this with the
protocol's literal \emph{unknown} rule, under which the scorable set is the
\TmpBothLitKnown{} $E$-allow cells and the agreement is
\TmpBothLitEeqY/\TmpBothLitKnown{}.

The stale model produces \TmpBothStaleFP{} false proofs (\TmpStaleFP{} on Gitea,
\TmpGStaleFP{} on GitLab), and a recrawl with the same credential removes
\TmpBothFixed{} of them, leaving \TmpBothStillMiss{} still missed and
\TmpBothLiveFP{} live false proofs (Table~\ref{tab:temporal}). The composition is
informative rather than uniform, and the two products agree on its shape: one
cell for the direct grant, one for the inherited-channel attachment, five and
six when a repository becomes public, and six and eight when a user becomes site
administrator. The last two are the cases a single-object probe is least likely
to anticipate, because the granting event is not about the object at all. The
false-proof phenomenon is also the part of the result that the
unknown-rule divergence does not touch: every stale false proof falls on an
$E$-allow cell, so both calibers of Table~\ref{tab:unknown-sensitivity} count
the same \TmpBothStaleFP{}, and the zero live false proofs likewise.

Revocations behave differently and must not be scored as misses. The old model
still lists \TmpBothOverRevoke{} cells whose access was revoked (\TmpOverRevoke{}
on Gitea, \TmpGOverRevoke{} on GitLab). In every one of them the request is
refused, so no violation occurred; treating them as false negatives would
require $E=\text{deny}$ \emph{and} $Y=\text{allow}$, which never happens here.

The multi-channel transition is where the products part company, and the
difference is itself a result. On Gitea, after the team member is removed,
\texttt{acme/t1} remains readable through a direct collaboration while
\texttt{acme/t2} becomes unreadable, and the ledger records the remaining
channel. On GitLab the same test has to be moved outside the group namespace:
removing a member from a group also deletes that user's project-level
membership in the group's projects---the project's direct-member list went from
one entry to none---so a direct grant placed inside the group would have measured
that product rule instead of multi-channel access. We therefore placed the
alternative channel on a personal project, and we record the cascade as a
product-semantics observation rather than encoding it as a ledger rule
(amendment A2 in the protocol record). This is the kind of non-local side effect
that makes a hand-maintained expectation model wrong in a way no single probe
can reveal.

For the role change, the read cell is unchanged while the write probe goes from
HTTP 403 with unchanged content to HTTP 201 with the owner re-reading the new
bytes; the owner then deletes the probe branch and observes 404, and the
repository's \texttt{README} digest is unchanged after cleanup.

The oracle column separates the products more sharply than the transition
counts. On Gitea the owner vantage answers \TmpOracleReadable{} of \TmpCells{}
cells (the remainder are anonymous, for which no oracle credential exists), but
the subject's own vantage answers only \TmpOracleSelfReadable{}---all of them
after the site-admin grant; outside it, a read collaborator receives HTTP 403 and
an outsider HTTP 404. On GitLab the owner vantage answers \TmpGOracleReadable{}
of \TmpGCells{}, and there the subject's own credential answers
\TmpGOracleSelfReadable{}: GitLab's member endpoint is readable by any project
member. That readability is not soundness. The endpoint answers ``is this
principal a member'', so it returns 404 for a public project a non-member may
read and for a site administrator who may read everything. Measured against the
ledger, the Gitea oracle disagrees on \TmpOracleDisagree{} cells and the GitLab
oracle on \TmpGOracleDisagree{}, all of the latter in the public
(\TmpGOracleDisagreePublic{}) or site-administrator (\TmpGOracleDisagreeAdmin{})
channel. A readable oracle can therefore be systematically wrong in exactly the
channels that are cheapest to overlook, which is why we never let it set $E$.
The two calibration regimes record this rather than hide it: \TmpCalChecks{}
known cells on Gitea, and on GitLab \TmpGCalGateChecks{} gated cells plus two
recorded diagnostics for the secondary vantage (amendment A1).

A channel-by-channel view makes the mechanism explicit
(Table~\ref{tab:oracle-channel}), for Gitea and GitLab. Each cell is assigned to
exactly one channel, and the rows sum to the cell count, so no cell is counted
twice; the grouping is by the strongest channel that grants it, not a
multi-label count. The same decomposition is available for Gogs
(Table~\ref{tab:gogs-channel}), and the reason an earlier draft omitted it is
worth recording, because it is a category error rather than a data limitation:
a channel label is a property of the entitlement ledger---built from the lab
topology, the documented inheritance rules and the recorded out-of-band grant
events, never from a probe response---so a product whose API exposes no
per-object membership route does not lose the channel axis. What such a product
loses is the ability to \emph{enumerate} those grants through that route, which
is a claim about the oracle surface, not about the labelling. Both tabulated products answer every
direct-collaborator and group cell correctly, and Gitea also answers its public
and site-administrator cells correctly. GitLab errs only by omission, and only on the two channels that
are not membership: it answers \emph{deny} on \TmpGOracleDisagreePublic{} public
cells that a non-member may read and on \TmpGOracleDisagreeAdmin{}
site-administrator cells. Neither product returns a single false allow. A
one-sided error is nevertheless sufficient to manufacture a false proof, because
an alarm on such a cell is then adjudicated against a deny that the ledger does
not support---which is the cross-vendor failure reported in
Section~\ref{sec:gogs}, now located in a named channel rather than attributed to
``the oracle'' as a whole.

\input{temporal_channel_table.tex}

Two properties of the run are recorded because they are what an external check can
verify. The protocol was frozen before the first Gitea transition, and its digest
is carried in both run records. For the amendments we establish the ordering by
content binding rather than by a self-reported clock: the archived pilot records
only A1, and the reported GitLab run records the digests of both A1 and A2, so
A2's content was in place when that run executed. The pilot is archived with its
cells excluded from every denominator. One recorded field does not support the
stronger reading and we do not use it: A2's \texttt{written\_at} says 02:32,
later than the run's last probe at 02:26:14, so the file's own timestamp cannot
establish that it was written before the first transition. We retain the field as
recorded and rely on the digest carried by the run; the discrepancy is documented
with the artifact rather than corrected in place. Inside a run, every cell carries its event sequence number, the
harness version before and after the transition, the $M_0$ and $M_1$ digests and
wall-clock timestamps, so the ordering---crawl, freeze, mutate, recrawl, oracle,
probe---can be re-checked line by line rather than trusted. Within the
enumeration, oracle and read probes, every call is a GET; the only mutating
calls are the transitions themselves, and
each is recorded with its status code. In addition to the read path, the write
capable transitions carry one write probe per cycle (an HTTP POST that creates
the probe artifact) followed by an owner-side DELETE that restores the state;
these write calls are recorded like every other call, they change no
authorization configuration, and they are the reason ``all observations are
GET'' would be too strong a summary of the run. Failed observations are not dropped
silently: a transport error after two retries is recorded as such and excluded
from the scored set, a mutation outside 2xx excludes its whole transition, and an
unreadable oracle is recorded as \emph{unknown} with the cell still scored on $E$
and $Y$.

Three scope limits are part of the result. First, the verdict is anchored to the
request outcome, and the transitions are serialised: no mutation is issued
between a cell's request and its probe, so the window between them is static by
construction. A deployment in which a grant or revoke lands between the request
and the adjudication is outside what these cells measure; for such a case the
cell's conclusion does not apply and the observation has to be repeated, and we
do not claim that the harness would detect the interleaving. Second,
\TmpBothStillMiss{} residual is
a count on these surfaces, not a completeness claim: the crawl sources
themselves define what can be seen, deploy keys are outside the ledger, and we
did not construct a channel that the recrawl is known to miss. Third, $hv$ is a
harness-maintained logical version, not an atomic server snapshot; the harness is
single-threaded and no concurrent writer ran, which is what makes the still
window meaningful. Fourth, the eight transitions are a designed set, not a
sample of deployment changes, so none of the counts is a field rate.

\input{temporal_table.tex}

\input{temporal_sensitivity_table.tex}

\subsection{Cross-vendor replication of the snapshot design: Gogs \GgVersion{} and GitLab \GlVersion{}}
\label{sec:gogs}

We repeated the same three-phase protocol and \ResRsCellsSnapshot{}-cell design on two further systems
(Table~\ref{tab:gogs-cross}): \textbf{Gogs} \GgVersion{}, the project Gitea
forked from in 2016 (related lineage, independently evolved code; census from its
pinned route table, \GgRoutePaths{} paths), and \textbf{GitLab} \GlVersion{}
(no lineage; census from the pinned running instance's routes, \GlRoutePaths{}
paths, since its OpenAPI document omits core grant endpoints). Gogs is a
replication across implementations, GitLab across vendors within one domain.

\input{gogs_cross_table}

The replication introduces a primary methodological variation: the
authorization-introspection surface, which shows that \emph{readability does not
imply soundness}. Gitea's owner-vantage
surface is readable and sound (agrees with effective access on all \GtOracleAgree{}
cells). Gogs's collaborator-membership endpoint is equally readable
(\GgOracleReadable{}/\GgCells{}) but unsound---wrong on \GgOracleWrong{} cells
(\GgOracleWrongPct\%), exactly the authorizations arising outside direct
membership: site administration (\GgFnSiteAdmin{}), public visibility
(\GgFnPublic{}), out-of-band grants (\GgFnOob{}). GitLab's
\texttt{members/all/\{u\}} resolves inherited membership, but public visibility
and site administration confer none: \GlOracleWrong{} false negatives
(\GlOracleWrongPct\%). What a surface cannot see is predictable from what it
measures---why clause~(2) of the admissibility condition
(Section~\ref{sec:admissibility}) cannot be automated from readability alone.

\input{gogs_channel_table.tex}

Table~\ref{tab:gogs-channel} locates those \GgOracleWrong{} errors instead of
leaving them as a single unsoundness figure, and it does so on the ledger axis
that the product's own API does not have to provide. Panel~(a) groups the same
\GgChCells{} cells by the granting channel recorded in the ledger. The surface
is correct on every cell it can represent: all \GgChCollaboratorAgree{} of the
\GgChCollaborator{} direct-collaborator cells and all \GgChNone{} cells with no
recorded grant. It answers \emph{deny} on all \GgChFalseDeny{} cells whose
entitlement arises from team membership (\GgChTeamFalseDeny{}), public
visibility (\GgChPublicFalseDeny{}) or site administration
(\GgChSiteAdminFalseDeny{}), and it never answers allow where the ledger denies
(\GgChFalseAllow{}). The error is therefore one-sided and located: exactly the
three mechanisms that the owner-vantage per-object collaboration route does not
represent. Panel~(b) re-groups the same cells by mechanism and shows the same
concentration as a rate rather than a count, which matters because the three
affected strata are small: \GgChSrcTeam{} team, \GgChSrcPublic{} visibility and
\GgChSrcSiteAdmin{} site-administrator cells, all wrong, against
\GgChSrcCollaborator{} and \GgChSrcNone{} cells, all right.

That the blind spot belongs to the chosen surface rather than to the product's
whole API is recorded, not inferred: in the surface probe shipped with the run,
the per-object permission suffix returns \GgChProbePermSuffix{}, the per-object
team route \GgChProbeTeamMember{} and the collaboration route the oracle
consults \GgChProbeCollabGet{}, while the organisation team listing is reachable
(\GgChProbeOrgTeams{}). A team listing is not a per-object authorization surface
and a surface-gated adjudicator does not consult it; the measurement therefore
reports what the gated surface licenses, not what the product could in principle
expose. The decomposition is also stable across readings: the live reading of
the same cells gives \GgChLiveAgree{} agreements and \GgChLiveFalseDeny{} false
denials, so the concentration is not an artifact of the snapshot reading.

Two labelling caveats are part of the table. First, \GgChMultiBasis{} cell
carries two ledger rules (site administrator and public visibility); it is
counted once, under the stronger rule, so the rows sum to \GgChCells{} cells
rather than to the \GgChRules{} recorded rules---the grouping is a partition of
cells, not of mechanisms. Second, the runner's provenance bucket is a heuristic
that files a collaborator grant made outside one particular repository as
out-of-band; the table therefore takes its labels from the ledger, and as a
cross-check the channel attribution of the \GgOracleWrong{} error cells agrees
between the two sources
(\GgChFnSiteAdminLedger{}/\GgChFnPublicLedger{}/\GgChFnTeamLedger{}).
The consequence is measurable: on Gitea the surface-gated adjudicators produced no
false proofs (\GtFpVthree{} v3, \GtFpDirectObserve{} direct-observe), while on
Gogs they manufacture them (\GgFpVthree{} snapshot, \GgFpVthreeLive{} live,
\GgFpDirectObserve{} direct-observe) and on GitLab again (\GlFpVthree{},
\GlFpDirectObserve{}), every v3 false proof falling on a cell whose authorization
derives from site administration or public visibility (\GgImplicitSiteAdmin{} and
\GgImplicitPublic{} such cells on Gogs; \GlImplicitSiteAdmin{} and
\GlImplicitPublic{} on GitLab). A readable membership surface with a blind spot converts that blind spot
into confirmed violations; no admissible oracle was available from the tested
introspection surfaces on Gogs or GitLab, so the honest adjudicator abstains on
the stratum where those surfaces would license a promotion.

Two results replicate everywhere: bookkeeping alarms false-proof on legitimately
granted access (Gitea \GtFpBookkeeping{}, Gogs \GgFpBookkeeping{}, GitLab
\GlFpBookkeeping{} cells)
with the same out-of-band-dominant grant profile; and the reachability mismatch
persists (Gitea \GtMismatchSnap{}, Gogs \GgMismatchSnap{}, GitLab
\GlMismatchSnap{} snapshot). On
Gogs the mismatch only partially resolves live (\GgMismatchLive{} remain---team
or public grants never enter a subject's personal crawl); on GitLab the live
reading recovers everything (\GlMismatchLive{} remain)---a second,
vendor-specific failure of the snapshot premise, in the crawl model rather than
the oracle.

\subsection{A seeded positive control}
\label{sec:seeded}

On the tested endpoints and transitions the three real systems enforce the
expected decisions, so the question remains: \emph{when a real object-level flaw
is present, does the same detector find it?} We answer with a positive control.
A \emph{controlled authorization target} faithfully emulates the Gitea API
contract the engine consumes, with a correct internal authorization model and
\SeedBolaN{} seeded vulnerable instances---all instantiating one fault pattern
(a none-entitlement subject served the object at
\texttt{GET /api/v1/repos/\{owner\}/\{repo\}}), differing in subject and object,
not in fault location; Table~\ref{tab:seeded} is therefore evidence about this
fault pattern, not three independent defect discoveries. Two design points. The
target deliberately exposes its authorization state to the subject-side harness
(mirroring the admissibility condition), which production Gitea does not---there
the oracle is readable from the owner/admin vantage
(\ResRsOwnerReadable{}/\ResRsOwnerReadableTotal{}) but not from an arbitrary
subject's (\ResRsBlackboxSelfReadable{}/\ResRsBlackboxSelfCells{})---and its
introspection surface reports the \emph{correct} \texttt{none} permission, so the
oracle stays sound and the defect is visible only as a discrepancy between oracle
and request. We
drove it with the same topology, mutation, crawl, introspection and adjudication
code as above, differing only in the system under test; ground truth comes from
the seed list, not from the detector.

\begin{table}[t]
\centering
\caption{Positive control: detection of \SeedBolaN{} seeded vulnerable instances
(one fault pattern). TP/FN/FP/TN are counted against the known seed, not against
a model.}
\label{tab:seeded}
\begin{tabular}{lrrrr}
\toprule
method & TP & FN & FP & TN \\
\midrule
\texttt{mr\_bookkeeping} & \SeedTPMrBookkeeping{} & \SeedFNMrBookkeeping{} & \SeedFPMrBookkeeping{} & \SeedTNMrBookkeeping{} \\
\texttt{v3} (owner vantage) & \SeedTPVThree{} & \SeedFNVThree{} & \SeedFPVThree{} & \SeedTNVThree{} \\
\texttt{v3} (subject vantage) & \SeedTPVThreeSelf{} & \SeedFNVThreeSelf{} & \SeedFPVThreeSelf{} & \SeedTNVThreeSelf{} \\
\texttt{direct\_only} & \SeedTPDirectOnly{} & \SeedFNDirectOnly{} & \SeedFPDirectOnly{} & \SeedTNDirectOnly{} \\
\texttt{direct\_observe} & \SeedTPDirectObserve{} & \SeedFNDirectObserve{} & \SeedFPDirectObserve{} & \SeedTNDirectObserve{} \\
\texttt{direct\_observe} (subject) & \SeedTPDirectObserveSelf{} & \SeedFNDirectObserveSelf{} & \SeedFPDirectObserveSelf{} & \SeedTNDirectObserveSelf{} \\
\bottomrule
\end{tabular}
\end{table}

Table~\ref{tab:seeded} is consistent with the negative result of the ablation in
this controlled fault pattern: \texttt{direct\_observe} and \texttt{v3} catch all
\SeedBolaN{} planted instances with zero false positives, while
\texttt{direct\_only}, which ignores the request outcome, raises
\SeedFPDirectOnly{} alarms on cells where the surface says \texttt{none} and no
access succeeds---alarms our strict definition does not promote. The control thus
narrows the claim to one fault pattern: when a surface that fully resolves the
authorization question is available, the metamorphic step is not necessary for
detection here. It does not generalise beyond that pattern, and it also delimits
the detector: we seeded one BFLA defect (non-privileged \SeedBflaActor{}
succeeding at \texttt{PUT \dots/collaborators/\{u\}}, HTTP \SeedBflaStatus{}),
which the read-oriented logic does not cover---a known boundary, reported as such
rather than as a detection failure. This is a controlled research artifact, not
production Gitea; a sanity check on the detector's logic, not a validation of
detection effectiveness on a population, and not a statement about Gitea's CVE
posture. Its source, seed list and run script are in the
reproducibility package at \url{https://github.com/ubfbuz3/JISA} (tag
\texttt{artifact-v1}).

\subsection{Four defect classes on the controlled target}
\label{sec:multiclass}

One fault pattern answers whether the detector finds the defect it was designed
for; it does not answer whether the chain transfers across authorization defects.
We therefore seed four classes on the same target, two independent instances
each, and drive them with the same crawl, introspection and adjudication code.
C1 is the read bypass already used. C2 is a write bypass: a subject whose true
permission is \texttt{none} changes the object's metadata with \texttt{PATCH}
and the change is observable afterwards. C3 is a faulty inheritance expansion: an
organization member outside the team is served the team's repository, and---this
is the point---the defect lives in the expansion itself, so the listing, the
search results and the object read all agree that the subject may see it. C4 is
an escalation: a subject that is neither owner nor administrator grants a third
party access it should not have, and the grantee can then reach the object; the
grant is recorded in the target's state but marked as obtained through the
defect, so the correct model still answers \texttt{none}. The classes are not
interchangeable: they stress different premises of the relation.

\input{multiclass_table.tex}

Table~\ref{tab:multiclass} reports per-class recall against the seeded truth. The
studied relation covers C1 and nothing else. C2 is invisible to it because the
relation is defined on whether the object is \emph{served}, not on whether it can
be \emph{changed}; moving the same relation to the channel the defect breaks
recovers both instances, which says something about the relation family rather
than about the relation we studied. C3 and C4 are worse than missed: in both the
subject already sees the object, so \texttt{cannotReachThroughGUI} is false and
the antecedent of the relation does not hold. The gate is not failing to notice
them; it reads its own premise and concludes there is nothing to test. That is a
property of reachability-gated testing against inheritance and escalation, not a
tuning problem, and it is the clearest evidence we have for why the gate cannot
be the whole adjudication.

The oracle-consuming adjudicator fires on all \McInstances{} instances once the
probe is channel-matched, and on \McOracleReadHits{} of \McInstances{} when
confined to the read channel. Its \McFpOracle{} false alarms on the
\McAuthorizedCells{} authorized cells are conditional on this target's oracle
being sound, and Section~\ref{sec:admissibility} reports what happens when it is
not. Two caveats belong here. First, the target's oracle shares the decision
function with the content path, so under the admission rule of
Section~\ref{sec:admissibility} every class is refused promotion: the
\McInstances{} hits are recall against known truth, not independence-certified
proofs. Second, these are four mechanisms on one emulator, not four discoveries
in production systems. Disabling any one class removes that class's alarms and
leaves every other class's alarms in place---each leave-one-out run returns zero
alarms on the disabled class---so an alarm is attributable to its own defect
rather than to a shared root cause.

\subsection{A publicly fixed defect}
\label{sec:cvepair}

The seeded control is a controlled artifact, so we also ran one defect that
nobody here planted: Gitea \CveVulnVersion{} against \CveFixedVersion{}, the
release whose notes record a security fix titled ``Fix comment permissions''
(\texttt{\#28213}, backported as \texttt{\#28216}); the same defect is catalogued
as CVE-2023-49946 for Forgejo, the fork that reported it. What is missing is
object ownership rather than path permission: a caller entitled to the
repository named in the URL can address an object that belongs to a different,
private repository.

Each version runs as a fresh instance, and the collision is built into the
topology: alice owns a private repository with one issue comment, bob owns a
public one. Bob sends \texttt{PATCH\allowbreak
/api\allowbreak/v1\allowbreak/repos\allowbreak/bob\allowbreak/pub\allowbreak/issues\allowbreak/comments/\allowbreak\{alice's comment id\}} with a replacement
body. On \CveVulnVersion{} the write succeeds (HTTP \CveVulnStatus{}) and
alice's own re-read returns the replacement body, so the effect is real and not
a status-code artifact. On \CveFixedVersion{} the same request is refused (HTTP
\CveFixedStatus{}) and alice's re-read is unchanged (Table~\ref{tab:cvepair}).
The ledger gives bob no write channel on alice's repository ($E$ = deny), so the
vulnerable cell is a true positive under the definition used throughout, and
the fixed cell raises no alarm: \CveTruePositives{} true positive and
\CveFalseProofs{} false proofs across the pair. All \CveControls{} control
writes behave as expected in both versions---bob's own comment is editable, the
correct path is refused, a third user is refused, anonymous is refused, an
absent identifier is refused---and alice's restore returns the body to the
original in \CveRestored{} of \CveVersions{} instances. The generating script is
fail-closed rather than trusting these fields: it requires the complete label
set with no duplicates, checks the recovery outcome, and recomputes $Y$ from
alice's re-read of the comment body before re-running the same
\texttt{classify}, refusing to emit a number if any layer disagrees.

What this establishes is narrow, and worth stating plainly. The relation is not
executed on this pair: the alarm recorded here is the $E/Y$ gate's own verdict,
so the result measures the gate, not the relations' recall on a real defect. It
shows that the gate can issue a correct violation claim and stops issuing it
once the vendor fixes the defect, so the gate is not merely a false-alarm
suppressor. It
does not show that we can \emph{find} such defects: the identifier is given, the
defect is public and already fixed, and the exercise is a positive control for
the gate, not a discovery. Nor is it an evaluation on a population---one defect,
one fault pattern (object versus path ownership), one product line. The
multi-defect coverage the same reviewer asked for is not obtained from this pair
but from the controlled target: Section~\ref{sec:multiclass} seeds read, modify,
inheritance and escalation faults and reports per-class recall, with the limits
of that construction stated there. The claim this pair supports is about the
gate's correctness on a real vulnerable/fixed pair, not about detection
coverage.

\begin{table}[t]
\centering
\caption{A publicly fixed defect (object-versus-path ownership). The verdict uses
$E$ from the ledger and $Y$ from the owner's re-read, not the write status.}
\label{tab:cvepair}
\small
\input{cve_pair_table.tex}
\end{table}


\section{Findings, admissibility, and falsification}
\label{sec:findings}

Each finding is a boundary on what the published relations can establish about
object-level authorization---\emph{temporal}, \emph{oracle} and
\emph{measurement}---carried by the studies of
Sections~\ref{sec:synthetic}--\ref{sec:realsystem}. The confirmations they rest
on (the predicates read at source, the published relations executed on a real
system, an independent implementation of the same predicates) introduce no
further claim and are collected as support evidence in Section~\ref{sec:support}.

\subsection{Finding 1 (RQ1): deterministic false proofs---the temporal boundary}
\label{sec:boundary1}

\emph{Taken as published---raw alarms, or alarms filtered by bookkeeping
exceptions---the adjudication reports a legitimate access granted out of band
after the crawl as an already-established violation, in the direction forced by
the two-valued predicate encoding.} On the constructed system, rate
\ResObsOobVtwoFpRate{} (\ResObsOobVtwoFalseProof{} of \ResObsOobFalseAlarm{}
alarms) against \ResObsFwVtwoFalseProof{} for the same share created in band; on
Gitea, \ResRsMrBookkeepingFp{} false proofs on \ResRsAuthorizedCells{} authorized
cells (\ResRsMrBookkeepingFpRate{}) under the crawl-time reading, against
\ResRsMrBookkeepingAlarmsLive{} alarms live---objects, entitlements and predicates
held fixed, only the reading time changed---so the error is a property of the
\emph{reading}, not of the system; the endogenous control
(Table~\ref{tab:rs-endo}) holds final visibility fixed and isolates that moment.
The direction is forced: the predicate has exactly two representations, and a
grant created outside the crawled interface has no representation but ``cannot
reach''. The predecessor attributes its residual false alarms to crawler
coverage~\cite{Mai2019Metamorphic}; we add that the failure is a function of
\emph{when} the model was read, which makes the natural oracle-based repair
available while showing no incremental adjudication benefit from the metamorphic
step (Finding~2).

\subsection{Finding 2 (RQ2): the metamorphic step adds no observed adjudication benefit}
\label{sec:boundary2}

\emph{If the false proofs of Finding~1 are removed by consulting a readable
introspection surface, it is the oracle that carries the work: for the
authorization relations and scenarios evaluated here, once the oracle fully
resolves the authorization proposition under test, replacing the entire
metamorphic step with a single ordinary request---one that consumes the same
readable oracle and access outcome---provides no observed incremental
adjudication benefit.} We state the
reason rather than present it as a discovered effect: its burden is an empirical
premise about the evaluated relations, not an algebraic identity.

We separate the two parts, because only one of them can fail. The algebraic part
is a lemma about function composition; the empirical part is a statement about
what these two relations actually read, and it has a falsifier.

\emph{The direct adjudicator.} Fix, before any run, the rule $D$ that issues
\texttt{violation-proven} for $(o,s,a)$ exactly when the oracle $O$ answers
\emph{deny} and the ordinary request $R$ succeeds. $D$ consumes $O$ and $R$ and
nothing else. It is a definition, not a result, and it is the comparator in the
ablation of Section~\ref{sec:synthetic-ablation} (\texttt{direct\_observe}).

\begin{lemma}[absorption]
\label{lem:absorption}
Let a detector issue its verdict through a map $g(A,O,R)$, where $A$ is its alarm
and $O,R$ as above. If, on the domain of cells considered, (i)~$A$ is a function
$a(M,R)$ of the crawl-time model and the response only, and (ii)~$M$ is
$O$-determined for the authorization proposition under test---that is, $M$ carries
no distinction the oracle does not already make---then there is a map $g'$ with
$g(a(M,R),O,R)=g'(O,R)$ on that domain.
\end{lemma}

\begin{proof}
By (i) the first argument of $g$ is a function of $(M,R)$; by (ii) $M$ is a
function of $O$ on the domain; substituting gives $A=a'(O,R)$, so $g$ factors
through $(O,R)$ as $g'(O,R):=g(a'(O,R),O,R)$.
\end{proof}

The lemma is deliberately weak: it says that \emph{if} a detector reads nothing
beyond the model and the response, its verdict cannot carry information the pair
$(O,R)$ lacks. It does not say that any particular implementation satisfies (i)
or (ii), and it does not compare two programs---it compares two functions on a
domain.

\begin{proposition}[the two relations fall in the lemma's domain]
\label{lem:oracle-suff}
For the two relations examined here, premise (i) holds by source inspection: their
predicates read the crawl-time reachability model and the response, and nothing
else (Section~\ref{sec:impl}). Premise (ii) holds on the cells tested, where the
oracle is readable and sound: $M$ and $O$ agree on the authorization proposition,
so $M$ adds no distinction. Consequently the verdict of each relation is
extensionally the direct adjudicator $D$ on those cells, and the equality is
checked, not assumed, by the paired ablation: the two methods disagree on
\ResAbDisagreementN{} of \ResAbNBola{} true-violation records
(Section~\ref{sec:synthetic-ablation}).
\end{proposition}

The proposition is falsifiable in one specific way: exhibit an observable $Z$
that a relation consumes, that is decision-relevant, and that $(O,R)$ does not
determine. Then premise (i) fails, the lemma does not apply, and the relation
could in principle beat $D$. We know of no such observable for these two
relations, and the ablation is consistent with its absence; that is the whole of
the claim. We do not assert that no such observable could exist for other
relations in the catalogue, and we do not assert that the relations are useless.

The proposition concerns \emph{adjudication} information only. It does not imply
that a metamorphic relation cannot improve input generation, search-space
reduction, or execution efficiency: a relation may guide which requests are
issued, or which gated subset is exercised, while carrying no verdict information
beyond $(O,R)$. It is a requirement on \emph{claims}, not a verdict on
metamorphic testing: it says only that a verdict computable from $(O,R)$ cannot be
shown to beat the direct consumer of $(O,R)$.

\paragraph{Scope, quantified}
The comparison above is pooled. Section~\ref{sec:stratum} re-runs it inside each
stratum of effective oracle readability (Table~\ref{tab:stratum}) and thereby
bounds the claim rather than
widening it. Where the oracle is readable (\MiiNOn{} records) the conclusion is
stronger than ``no observed benefit'': the plain request carries fewer false proofs
(\MiiOnDirectObserveFP{} vs.~\MiiOnVTwoFP{}) and no fewer detections. Where it is
not (\MiiNOff{} records) the antecedent fails---there is no oracle to consult, and
every oracle-bearing method detects nothing by construction---so the metamorphic
family is compared against oracle-free references instead: it retains an
advantage over a trivial floor and over a single uninterpreted request
(\MiiLiftVTwo{}$\times$ vs.~\MiiLiftObsSuccess{}$\times$ over the base rate), and
that advantage is exactly coextensive with what the crawl-derived model can see.
Finding~2 is a claim about the regime its antecedent describes; it does not say
the relations are worthless outside it.

Evidence: Table~\ref{tab:ablation}---\ResAbDirectObserveFP{} false proofs
vs.~\ResAbVThreeFP{} (a tie at zero, not a win) over the \ResAbNBola{}
true-violation records. The comparison is paired, so we report the discordance
rather than an inferential interval: the two methods disagree on
\ResAbDisagreementN{} records, and on every one \texttt{direct\_observe} is
right---the aggregate recalls \ResAbDirectObserveRec{} and \ResAbVThreeRec{}
restate that asymmetry rather than measure a difference, and no margin was fixed
in advance. Gitea independently reproduces the zero-false-proof result
(\ResRsDirectObserveFp{} vs.~\ResRsVThreeFp{}), but not the same comparison,
because the tested surface contains no seeded vulnerability; the cross-vendor
systems (Section~\ref{sec:gogs}) and the seeded control (Section~\ref{sec:seeded})
complete the evidence. One qualification is part of the claim: the control is
\emph{constructed}, not discovered---\texttt{direct\_observe} instantiates the
direct-consumer construction $D$ of Proposition~\ref{lem:oracle-suff}, so the tie
is the empirical face of the necessary condition, a statement about the
definitions rather than a coincidental measurement.

\subsection{Finding 3 (RQ3): the reported specificity is not conditioned on the correspondence it assumes}

\emph{The reported \ResLitReportedSpecPct{}\% is a statement about a corpus of
generated inputs; it is not conditioned on an explicit check that the
reachability model coincides with the true authorization state, and its
input-level denominator cannot, on its own, provide stable scenario-level
resolution for a scenario-uniform failure.} The omission:
adjudication is a function of the
reachability model, so where model and true state diverge---the out-of-band case
of Finding~1---the figure measures agreement with the tool's own crawler; the
authors concede the mapping is ``not always
available''~\cite{Chaleshtari2022Metamorphic}, but the number is reported
without the condition. The denominator: a failure uniform across a scenario
class contributes only as many false positives as that class emits inputs, and
is diluted by the rest of the run (Proposition~\ref{prop:resolution}: in our
approximation the configuration that misreports a whole class of legitimate
access still scores \ResSpecRowOneSpecPct{}\%; Corollary~\ref{cor:resolution}:
no fixed threshold can guarantee that such a class forces the metric below it).
Nor does the report fix the weights: if a failing class supplies a share $w$ of the
follow-up inputs, the figure moves only between $1$ and $1-w$, so telling a
small-share class failure apart from no failure at all requires an estimate of $w$
that the published number does not carry.
We are not saying
the figure is false or the false positives misdiagnosed; the claim is that such a
metric does not itself validate the reachability-to-authorization correspondence
on which its interpretation depends, and that a class-uniform failure can be
held above any fixed threshold by unaffected inputs.

\subsection{Admissibility, and falsification}
\label{sec:admissibility}

Three properties must be kept apart. A surface is \emph{readable} when its answer
can be obtained from the vantage in question; \emph{sound} when its answer equals
the true authorization state; and \emph{admissible}, for promotion to
\texttt{violation-proven}, when (1)~it is readable \emph{and} (2)~one can give a
reason that it is not of the same kind as the defect under test---not itself a
reachability artifact whose blind spot correlates with the blindness tested.
Readability is what makes soundness usable, but it neither implies soundness nor
confers admissibility: the owner-vantage authorization-introspection surface is
readable on every cell in each system (\ResRsCellsSnapshot{} Gitea, \GgCells{}
Gogs, \GlCells{} GitLab), yet sound on Gitea alone---it agrees with effective
access on \GtOracleAgree{}, \GgOracleAgree{} and \GlOracleAgree{} of those
cells respectively. The
cross-vendor systems are therefore not cases of ``no oracle'': a readable
membership surface exists on both, but it is unsound for the authorization
channels exercised here, so it fails clause~(2). Clause~(2) is a per-system
judgement and not automatable from readability alone; the cross-vendor results
make it concrete (Gogs and GitLab fail it; their adjudicators convert the blind
spots into \GgFpVthree{} and \GlFpVthree{} false proofs). Where the condition
holds, the honest method consumes the oracle directly and drops the relation;
where it does not, the honest method abstains---as v3 does on the unreadable
stratum, where the aggregate table would have ranked it best, a ranking that
reflects abstention rather than improved detection.

Read together, the two findings bound what the technique's adjudication can yield
for the authorization proposition studied. If the reachability stand-in is the
only warrant and it is wrong, a legitimate access becomes a false proof
(Finding~1). If a readable, sound oracle is instead available, the relation adds
no adjudication over the direct consumer of that oracle and the outcome
(Finding~2). And where no readable, admissible oracle exists---the black-box
case---there is nothing for a sound method to consume, so it abstains rather than
proves. Across the statuses the oracle can take, the verdict the evaluated
relations can yield is thus a false proof, a result no better than the plain
request (and, where the oracle is readable, worse), or an abstention. None of
those is a proof of illegal access, because the relation's alarm is not evidence
of it: an alarm says that the relation failed, and an ordinary request that
succeeds says that the object was served. A genuine violation becomes reportable
only when an independent entitlement source says \emph{deny} and a content-level
effect observation says \emph{allow}; Section~\ref{sec:cvepair} reports that gate
producing one true positive on a vulnerable build and none on its fixed
successor. The relation contributes to neither side of that conjunction.

Each finding would fall as
follows: Finding~1, if the predicates were defined over live rather than
recorded reachability; Finding~2, by an authorization MR needing an observable
that a succeeding request plus a readable oracle cannot produce; Finding~3, by
a specificity in scenario units, or with the condition explicit, that still fails
to separate the configurations studied here.

Table~\ref{tab:evidence-claim} states, for each research question, what is
assumed, what was measured, and what the result does not cover. We add it because
the three findings rest on different kinds of evidence---source-level functional
dependence, a constructed enumeration, and two real-system studies---and a reader
should be able to see which kind carries which sentence.

\begin{table*}[t]
\centering
\caption{Evidence to claim. Each row names the assumption the claim needs, the
artefact that carries it, and the boundary that remains. The last row is not a
research question: it checks the adjudication gate used throughout, and is not
evidence about detection effectiveness.}
\label{tab:evidence-claim}
\footnotesize
\begin{tabular}{@{}p{0.055\textwidth}p{0.235\textwidth}p{0.185\textwidth}p{0.245\textwidth}p{0.205\textwidth}@{}}
\toprule
RQ & Claim & Assumption it needs & Evidence & What it does not cover \\
\midrule
1 &
A crawl-time reading reports an out-of-band grant of legitimate access as a
proven violation. &
The grant occurs after the crawl, is recorded in the ledger, and is reachable
through a channel the crawl does not enumerate. &
\ResSynRecords{}-record constructed study; \TmpBothPolicyEvents{} committed
transitions on Gitea and GitLab, \TmpBothStaleFP{} stale false proofs, all
\TmpBothFixed{} removed by a same-credential recrawl. &
Two relations of one reachability-gated family; designed transitions, not a
sample of deployment changes; no residual miss is a completeness claim. \\
\addlinespace
2 &
Once a readable, sound oracle is supplied, the relation adds no observed
adjudication benefit over the direct consumer of the same oracle and outcome. &
The oracle is readable and sound on the cells scored; the relation consumes no
decision-relevant observable outside $(O,R)$. &
Paired ablation on \ResAbNBola{} true-violation records, disagreeing on
\ResAbDisagreementN{}; channel-level oracle fidelity over
\TmpBothCells{} cells. &
Conditional on the antecedent: the unreadable stratum has no oracle to consult,
and the oracle shares its implementation with enforcement. \\
\addlinespace
3 &
The published specificity is an input-level figure whose denominator cannot, on
its own, expose a class-uniform failure. &
Scenario classes emit comparable numbers of inputs, so a uniform failure is
diluted by unaffected inputs; the class shares are not reported, so the
dilution is bounded rather than estimated. &
Proposition~\ref{prop:resolution} and its corollary; sensitivity of the example
to the assumed inputs per case. &
The original per-case denominators are unavailable, so the published figure is
diagnosed rather than re-estimated. \\
\addlinespace
Gate &
An independent entitlement ledger plus a content-level effect observation
reports a real defect as a true positive and stays silent on its fix. &
The ledger is complete for the tested object and the effect is read back by the
owner, not inferred from a status code. &
Gitea \CveVulnVersion{} versus \CveFixedVersion{}: \CveTruePositives{} true
positive, \CveFalseProofs{} false proofs, \CveFixedAlarms{} alarms on the fixed
build, \CveControls{} control probes as expected. &
One publicly fixed defect found and supplied by us; this is a check of the gate,
not a recall estimate. \\
\bottomrule
\end{tabular}
\end{table*}

\section{Support evidence}
\label{sec:support}

Nothing here is a further finding. This section holds what the three findings
rest on: the semantics of the two governing predicates read at source
(Section~\ref{sec:impl}), an execution of the published relations on a real
system (Section~\ref{sec:native-run}), and the robustness of our numbers to a
second, independent implementation of the same predicates
(Section~\ref{sec:secondimpl}); the seeded positive control
(Section~\ref{sec:seeded}) and the cross-vendor replication (Section~\ref{sec:gogs}) are the
remaining confirmations and stay with the study that produced them. Each item is
a confirmation of a mechanism already stated, not a new claim about what the
technique can detect.

\subsection{Confirmation at source}
\label{sec:impl}

Definitions can be read generously, so we checked the predicates at source in
the public library~\cite{Chaleshtari2022Metamorphic}. It ships \ResArtMrTotal{}
relations, \ResArtAuthzMrTotal{} authorization relations,
\ResArtAuthzGatedEither{} gating on the two predicates analyzed. Four
implementation-level properties not made explicit in the main presentation:
(1)~\emph{the freeze is hard-coded}---%
\texttt{cannotReachThroughGUI} is set-membership over the crawl inputs, issues
no request, and is cached in a field written once and never invalidated
(\texttt{userCanRetrieveContent} likewise reads a crawl-produced store loaded
once); (2)~\emph{the clause that could have caught it is inert}---the
``absent across all users'' set is built by stopping each recorded sequence at
the first login action, so it is always empty; (3)~\emph{missing data defaults
to an alarm}---with an empty or null crawl argument the predicate returns
``cannot reach'', so an empty crawl \emph{satisfies} the precondition: the
mechanism behind Finding~1's deterministic direction; (4)~\emph{the role
relation is configuration, not state}---the supervisor guard reads a static
configuration file and reports true for identical accounts, so it cannot reflect
a later grant.

\subsection{Executing the published relations}
\label{sec:native-run}

The library ships relations but no execution environment; we supplied the system
configuration and crawl snapshot and ran the relations---source unmodified,
build configuration (\texttt{pom.xml}) and an external harness added
(\NtUpstreamModified{} changed upstream file, \NtUpstreamJavaModified{} changed
Java sources, two added classes totalling \NtHarnessLines{} lines)---against
Gitea \NtGiteaVersion{} (Table~\ref{tab:native-facts}). \NtMrExecuted{} of \NtMrTotal{} authorization relations
executed; \NtMrFired{} returned failures, \NtAlarms{} alarms in all---the
engine's own \texttt{getFailures()} value (Table~\ref{tab:native}). What the run
establishes is scoped to the \emph{reachability-gated subset that fires here},
not a full replay of the catalogue: one relation aborts on an engine-side defect,
most issue no request because their own guards are unsatisfied, and the content
predicate is never invoked. Almost none
of the \NtElapsedS{} seconds is analysis: each leaf evaluation invokes a
browser-cleanup hook costing minutes, a cost of the missing environment that
bounds any reproduction.

\begin{table}[t]
\centering
\caption{Execution facts for the native run of the published relations
(Section~\ref{sec:native-run}). All values are emitted by
\texttt{ingest\_native\_engine.py} from the run record; upstream Java sources are
unchanged apart from the build configuration.}
\label{tab:native-facts}
\small
\begin{tabular}{@{}p{0.66\columnwidth}r@{}}
\toprule
item & value \\
\midrule
authorization relations: catalogue, executed, fired & \NtMrTotal{}, \NtMrExecuted{}, \NtMrFired{} \\
engine alarms (the engine's own return value) & \NtAlarms{} \\
HTTP requests attributed to observing relations & \NtHttpCalls{} \\
changed upstream Java sources & \NtUpstreamJavaModified{} \\
build/config changes (one \texttt{pom.xml} line) & \NtUpstreamBuildModified{} \\
added harness classes / added lines & two / \NtHarnessLines{} \\
\bottomrule
\end{tabular}
\end{table}

The load-bearing cell: the subject \NtKeyGrantee{} holds a legitimate
entitlement to \NtKeyRepo{}, granted out of band after the crawl. There the
published \texttt{cannotReachThroughGUI} answers ``cannot reach''---the answer
it is \emph{defined} to give, since the subject's crawl record has no entry for
the repository; across the run it answered so \NtCrtgTrue{} of \NtCrtgCalls{}
invocations (\NtCrtgTruePct{}\%), \NtKeyCrtgTrue{} on this cell, and
\NtKeyMrFired{} relations (\NtKeyMrNames{}) report a failure there. The access
is authorized: all \NtByteCells{} granted cells were re-probed, and the
grantee's response hashes identically to the owner's while a third account is
refused. The \NtAlarms{} alarms fall on \NtAlarmCells{} cells, all post-crawl
grants\ifnum\NtAlarmsOffGrant=0, none on a non-entitled object.\else, plus
\NtAlarmsOffGrant{} on a non-entitled object.\fi{} Limits: the harness's HTTP
observation channel is exact for the reachability predicate (set-membership, no
request) but unexercised for the content predicate (\NtUcrcCalls{} invocations;
supervisor guard \NtSupCalls{}, \NtSupTrue{} true); not all gated relations fire,
for reasons inside their own guards, and \NtMrErrored{} abort on an engine-side
defect---so the reproduced false proof is that of the reachability-gated subset
that fires, not a full-replay count of the published relations. And Gitea, on
the tested surface, enforces the expected decisions, so the run can only exhibit
false proofs.

\paragraph{Which stage of that claim does the native run establish?}
Three stages have to be kept apart, and only the third is a security finding.
The \emph{native alarm} $A$ is what the engine itself returns:
\texttt{getFailures()} is appended by \texttt{MR.fail()} (\texttt{MR.java}:909),
which logs \texttt{FAILURE} together with the failing input pair and asserts
nothing about authorization---no ledger, no expected decision, no verdict.
Reading a fired authorization relation as a \emph{violation claim} $H$ is a
promotion policy $P$ that we state rather than inherit: $P$ maps every fired
\texttt{OTG\_AUTHZ} relation to ``a violation was observed'', which is the
reading the catalogue's purpose invites but its code does not perform. Since $P$
applies to every alarm, $H=\NtStageH$ holds by construction and carries no
empirical content. All of it is in the independent truth $T$, computed from the
ledger grant ($E$) and a content re-read ($Y$: the responder's bytes are
identical to the owner's and a third account is refused): \NtStageTTrue{}
violations are confirmed and \NtStageTFalse{} of the \NtStageA{} promoted alarms
are contradicted, on \NtStageTCells{} cells (Table~\ref{tab:native-split}). A
reader who declines $P$ should read the run as \NtStageA{} failed metamorphic
relations rather than as \NtStageA{} vulnerabilities; the disagreement survives
that reading, because the failures fall on accesses the ledger allows.

\begin{table}[t]
\centering
\caption{The native alarms split into the three stages: $A$ (the engine's own
\texttt{getFailures()}), $H$ (the promotion policy $P$ applied to $A$), and $T$
(independent truth from the ledger $E$ and the content re-read $Y$). Every row
is an alarm, so $A$ is not repeated as a column; $H$ likewise applies to every
row. $^{\dagger}$responder bytes identical to the owner's (sha256) and a third
account refused with 404. MR names abbreviate \texttt{OTG\_AUTHZ\_}$x$.}
\label{tab:native-split}
\footnotesize
\input{native_split_table.tex}
\end{table}


\begin{table*}[t]
\centering
\caption{The published authorization relations executed against Gitea
\NtGiteaVersion{}. Each row is the applicability chain from the catalogue to the
outcome: \emph{status} is the engine's per-relation result (\NtMrExecuted{} of
\NtMrTotal{} evaluated, \NtMrErrored{} aborted on an engine-side defect),
\emph{alarms} its own \texttt{getFailures()} return (\NtMrFired{} relations,
\NtAlarms{} in all), and \emph{requests} (\NtHttpCalls{} total) attributed to the
relation observing them first.}
\label{tab:native}
\small
\input{native_engine_table.tex}
\end{table*}

\subsection{Implementation robustness}
\label{sec:secondimpl}

An independently written implementation of the same predicates was compared with
the main one on a subset fixed before the comparison, not sampled after it: the
in-band scenarios on which a second reading is defined, holding the
exception-policy variant at the representative level and the reachability reading
at the one each relation uses (MR-002 blind, MR-004 aware). This yields
\ResAltPoints{} (configuration, scenario) points; MR-002 is defined on
\ResAltMrTwoPoints{} of the resulting relation comparisons and MR-004 on the
remaining \ResAltMrFourPoints{}, giving \ResAltCompared{} in all. The two agreed
on \ResAltAgree{} of them (\ResAltAgreeRate{}); the point set is a fixed
enumeration, so no interval is attached. Both were written from the same
reading of the semantics, so this bounds implementation accident only; it does
not bound a shared blind spot, and we do not present it as corroboration.

\section{Threats to validity}
\label{sec:threats}

\paragraph{Source versus execution} We read the predicates at source (built
under JDK~8 and Maven; Section~\ref{sec:impl}) and then executed the relations,
unmodified apart from the build configuration and our harness
(Section~\ref{sec:native-run}). The harness substitutes an HTTP observation
channel---exact for the reachability predicate, not exercised for the content
predicate, whose invocation count we report rather than claim cleanliness we did
not demonstrate. The run converts one class of claim (the false proof is produced
by the \emph{published} relations on a real system); it does not convert the
rest: the scenario/configuration and cell studies remain statements about our
instantiation,
our exception vocabulary belongs to the same predicate family, and we report no
sensitivity or specificity attributable to \mbox{MST-wi}.

\paragraph{Metric and statistics} The specificity approximation uses two
follow-up inputs per case---not the engine's count; the conclusion about the
denominator's resolution does not depend on the exact value. Records within a
design cell are not independent (Section~\ref{sec:protocol}, statistical
treatment) and the records are a
fixed enumeration, so we report counts and rates rather than inferential
intervals, and no result rests on statistical significance. The RQ2 ablation is
paired and is reported as a discordance count.

\paragraph{Two MRs, one family} We instantiated \ResDesignMrs{} MRs of one
family---relations gated on the absence of evidence. The native run lifts the
count to \NtMrExecuted{} of \ResArtAuthzMrTotal{} catalog authorization
relations but not the family: only the reachability-gated ones alarm in our run.
This is the paper's single most important limitation; raising it requires
relations gated otherwise, a different experiment.

\paragraph{Real systems} Three systems (one lineage pair, one independent),
chosen for accessibility, none vulnerable on the tested surface---so only the
false-proof side is measured, within one domain, the git forge. Channel
proportions are topology design parameters, not field prevalence. Constructed
scenarios with construction-defined ground truth are why the real-system section
exists; three systems tested once each are still not a population.

\paragraph{Stratification and the references used in it} The split of
Section~\ref{sec:stratum} is on whether the introspection surface returned a
reading for the record, not on the configuration switch: \MiiFlagOnButUnavailable{}
records sit in configurations that enable introspection but whose probe returns
nothing, so the switch would misclassify them. The two oracle-free references
there are deliberately weak---one unconditional, one guessing from a single
request---because the point is to size the metamorphic step against methods that
need nothing, not to identify the best oracle-free detector; a differential
prober across accounts consumes more than those references do and was not
instantiated. Both references share an observable with the labelling rule, as
\texttt{direct\_observe} does, so only the differences against them---not their
absolute rates---are meant to be read.

\paragraph{Our own earlier claims} An earlier iteration proposed v3 as a
mechanism on the strength of its precision (\ResAbVThreePrec{}), which was
bought with \ResAbVThreeFN{} false negatives, and a control
(\texttt{direct\_only}, precision \ResAbDirectOnlyPrec{}) that discarded the
access result. Replacing it with a fair control produced the negative result. We
record this as the condition any similar claim must meet: a control consuming the
same observable quantity as the ground truth.

\section{Related work}
\label{sec:related}

\paragraph{Metamorphic security testing} Metamorphic testing was brought to web
security by Mai et al.~\cite{Mai2019Metamorphic}, who cast security test
patterns---authentication, authorization, session---as metamorphic relations so
that faults can be exposed without an oracle. The line matured into
\mbox{MST-wi}~\cite{Chaleshtari2022Metamorphic}, whose \ResLitMrCatalog{}
system-agnostic relations over a Web-specific domain language are the object of
our study. We take the
technique's own admission that the role-to-resource mapping is ``not always
available'' at face value and follow it to its consequence for what the
adjudication can decide. A relation whose guard reads a crawl-time model is
sound only while the model equals the authorization state, and that equality is
exactly what the technique cannot check for itself.

\paragraph{Detecting broken access control} A large and active literature
detects access-control flaws directly, and it is worth separating the oracles it
relies on. Black-box probing drives requests with and without the rights of
another subject and looks for a missing refusal~\cite{Liu2025BACScan,Dharmaadi2025BACFuzz},
a line that runs back to AUTHSCOPE~\cite{Zuo2017AUTHSCOPE} and that has been
extended to deployed systems under ethical constraints~\cite{Chehade2025Forbidden};
where an API description is available, the oracle itself can be synthesised, for
instance by transforming an OpenAPI specification into a colored Petri net that
predicts per-object access~\cite{Filho2025Automated}.
A complementary, white-box line recovers the \emph{enforced} policy from the
implementation---by reverse-engineering RBAC policies~\cite{Le2021Automated}, by
statically extracting authorization policies~\cite{Berger2020Static} or database
access-control rules~\cite{Zhang2026Extracting}, or by locating the permission
checks themselves~\cite{Shi2025Facilitating}---and then flags violations against
that recovered model: BOLA in database-backed applications~\cite{Huang2024Detecting},
missing-owner-check and function-level flaws in Java web
applications~\cite{Liu2025MOCGuard,Feng2025JAuthGuard}, function- and object-level
access-control flaws across Web applications~\cite{Schlaubitz2025A}, and privilege
escalation in polyglot microservices~\cite{Li2026Detecting}. Specification-driven API
fuzzers generate the requests well but inherit the same oracle
difficulty~\cite{Atlidakis2019RESTler,Arcuri2024Advanced,Seran2025Handling};
adding explicit access-policy violation checks~\cite{Sahin2027Enhancing} moves
them toward the detection problem but does not remove the need for a ground
truth, and a recent turn delegates generation or probing to language
models~\cite{Uddin2026Automated,Sun2025LLM,Wang2025Large}. A concurrent line
reframes BOLA under a zero-trust model, in which each request must carry an
explicit authorization decision rather than one inferred from endpoint
reachability~\cite{Wu2026Rethinking}; our question sits at that decision level.
Each of these papers reports detections against an oracle it holds
privately---a policy, a seeded defect, a differential response---and none of them
is our comparison class, because we do not measure detection rates. What this
paper adds is three claims about one published adjudication, each with a
falsifier. Two of the ingredients are already known---that a crawl-time model
can go stale, and that an authorization oracle can be unsound; the claims below
differ in that the first failure is deterministic and attributable to a
representation choice, the second is bounded by an explicit equivalence rather
than argued case by case, and the third is a property of the published metric's
denominator. First, a gate that reads a crawl-time model misreports a grant the
record does not contain, deterministically rather than intermittently
(Finding~1; the constructed enumeration). Re-reading the surface repairs that
misreport only when the grant is one the surface would have shown on a second
reading: across the two products a same-credential recrawl removes
\TmpBothFixed{} of \TmpBothStaleFP{} stale false proofs, leaving
\TmpBothLiveFP{} live ones. When the grant arrived through a channel the
surface never exposes, no re-reading of that surface repairs it---on Gogs
\GgMismatchLive{} of \GgMismatchSnap{} mismatching cells and \GgFpVthreeLive{}
of \GgFpVthree{} surface-gated false proofs persist under live reading, all on
public-visibility or team-derived grants that never enter a subject's personal
crawl. Staleness is repairable and channel blindness is not, and the two call
for different remedies. Second, once a readable,
sound oracle is supplied, the promoted verdict is a subset of direct adjudication
on the evaluated domain, so the metamorphic step adds no confirmation there
(Proposition~\ref{lem:oracle-suff}(b)); the falsifier is a record the relation
confirms and the direct consumer does not. Third, the published specificity's
denominator is pooled over the inputs a crawl could enumerate, so a failure
confined to a class outside that set does not move it
(Proposition~\ref{prop:resolution}). We claim no new detector, no new oracle,
and no recall advantage over the access-control line; what we claim is a
statement about what an oracle-substituting technique can establish, and where
that statement stops.

\paragraph{Direct authorization oracles, and what is at stake here} Comparing
responses across principals or privilege states is an established way to obtain
an authorization oracle when the ground-truth relation is available, and the
black-box and white-box lines above are instances of it. We propose no new
direct oracle, and we do not claim that direct authorization comparison is
novel. Our subject is a different question: what an existing
oracle-substituting technique \emph{decides} when it substitutes reachability
for entitlement. The direct oracle enters our design in exactly one
role---as the fair, same-observable control that a necessity claim about the
metamorphic step must face. Without such a control, a substitution can appear
necessary only because the alternative it is compared against was not allowed to
consume the same decision-relevant observables.

\paragraph{LLMs and the oracle problem} LLMs are used for vulnerability detection
and specification inference~\cite{Wang2024A,Wen2024Enchanting,Wang2026Measuring,LeCong2025Can},
including lightweight models over hybrid code features~\cite{Liu2024Enhancing} and
federated, explainable detectors~\cite{Senanayake2024Defendroid}; we use none in the
detection loop. The transferable condition this paper exposes
is that any claim of necessity must be tested against a fair substitute that
consumes the same observable quantity as the ground truth.

\paragraph{Measurement and evaluation} Our contribution is finally a measurement,
and it connects to work on the quality of the artifacts security studies are
built on---for instance the observation that vulnerability datasets are
frequently mislabeled, duplicated or otherwise unfit for the claims drawn from
them~\cite{Croft2023Data}. The artifact we report is of the same family but
inside a study rather than a dataset: a specificity figure whose
\emph{denominator} is pooled over the inputs a crawl could enumerate, so it
cannot register a failure confined to a class that crawl never reached, while
the adjudication that consumes the same crawl can. We therefore read a testing technique's own reported metric
as a measurement instrument, stating explicitly what its denominator can and
cannot resolve.

\section{Conclusion}
\label{sec:conclusion}

For object-level authorization, \mbox{MST-wi}'s predicates do not know whether
the acting subject is entitled to the object; they know what the interface, as
crawled, permitted. Instantiating them as defined produced three findings
(Section~\ref{sec:findings}).

\paragraph{Finding 1 (temporal)} As designed, the adjudication false-proofs on
legitimately granted access, deterministically (\ResObsOobVtwoFpRate{}
constructed; \ResRsMrBookkeepingFpRate{} Gitea), because an out-of-band grant fits
only the wrong one of the predicate's two representations. On the real system the
mismatch appears only under the crawl-time reading
(\ResRsMismatchSnapshot{}/\ResRsMismatchDenom{}, vs.\ \ResRsMismatchLive{} live),
with an endogenous control holding final visibility fixed: within this controlled
design the only manipulated explanatory variable is when the model was read, not
what was visible.

\paragraph{Finding 2 (oracle)} The natural repair---a readable, sound
authorization oracle that fully resolves the proposition under test---works, and
in that regime the evaluated metamorphic step
provides no observed incremental adjudication benefit over a direct consumer of the same oracle and outcome:
an empirical premise, formalised as a conditional (Proposition~\ref{lem:oracle-suff}), whose empirical face is that
a single ordinary request yields the same false-proof count (false proofs
\ResAbDirectObserveFP{} vs.~\ResAbVThreeFP{}), and that of the \ResAbNBola{}
true-violation records the paired comparison disagrees on only
\ResAbDisagreementN{}, each time in favor of the plain request; no margin was
pre-specified, so no inferential claim attaches.
On Gitea the comparison is limited to false-proof behavior, because no
vulnerability was present on the tested surface.

\paragraph{Finding 3 (measurement)} The reported specificity of
\ResLitReportedSpecPct{}\% is diagnosed as a specific published figure: it is
$\mathrm{TN}/(\mathrm{TN}+\mathrm{FP})$ over the follow-up inputs the crawl
could enumerate, with manually verified labels, and that denominator cannot
register a failure confined to a class the run did not exercise.

Support evidence: reading the implementation confirmed the semantics and added
the mechanism---the crawl model is frozen by a cache never invalidated, and a
missing crawl defaults to ``unreachable''---and the positive control is a sanity check that closes the loop: on a seeded
vulnerable target the surface-gated adjudicator detects all planted instances of
the fault pattern it was built for with zero false positives, and on four seeded
defect classes it detects the read bypass and nothing else---the write bypass
breaks a channel the relation does not test, and the inheritance and escalation
faults falsify the antecedent because the subject can already see the object.

\paragraph{What the two added experiments carry} The temporal finding is not
confined to one product. Under the same pre-registered protocol, frozen before
the second run was observed, GitLab CE \GlVersion{} commits \TmpGPolicyEvents{}
transitions over \TmpGCells{} cells, with entitlement and effect agreeing on
\TmpGEeqY{} of them: \TmpGStaleFP{} stale false proofs, all \TmpGFixed{} removed
by a recrawl with the same credentials, \TmpGLiveFP{} live false proofs and
\TmpGStillMiss{} residual misses. Pooled with Gitea this is \TmpBothCells{}
cells, \TmpBothEeqY{} agreements, \TmpBothStaleFP{} stale false proofs and
\TmpBothFixed{} fixes. The replication is of a protocol, not of a topology:
GitLab's project-membership surface answers 200 or 404, so it errs only by
omission---in \TmpGOracleDisagree{} cells it disagrees with the ledger, every one
of them a false deny, concentrated in the public
(\TmpGOracleDisagreePublic{}) and site-administration
(\TmpGOracleDisagreeAdmin{}) channels---and both properties are read from the
same run rather than one being derived from the other. The entitlement/effect
gate is in turn checked against a real defect: on a publicly fixed
object-versus-path ownership flaw, Gitea \CveVulnVersion{} yields
\CveTruePositives{} true positive with \CveFalseProofs{} false proofs, and the
fixed \CveFixedVersion{} yields \CveFixedAlarms{} alarms over \CveControls{}
control probes. That is a check of the gate on one defect we located and
supplied; it is not a detection rate.

\paragraph{What we do not claim, and how far the claims reach} We do not claim
that \mbox{MST-wi} is wrong or say anything about its accuracy as a detector; the
native run establishes a qualitative outcome, for one gate family, on one real
system that on the tested endpoints and transitions enforces the expected
decisions. Three ceilings stack and bound everything above. First, our two
relations are two of a catalogue of \ResLitMrCatalog{}, and both belong to one
reachability-gated family. Second, even the execution of the published relations
on the native engine reaches only that family: \NtMrFired{} of the \NtMrTotal{}
authorization relations fire (\NtMrExecuted{} executed; \NtKeyMrNames{}), and the content predicate
\texttt{userCanRetrieveContent} is never invoked (\NtUcrcCalls{} calls), so that
evidence is bounded the same way. Third, the real-system study exercises the
\emph{false-proof} face on a single constructed fault pattern, not the detection
face against a population of real flaws, and the Gitea channel proportions are
topology design parameters, not field prevalence. Together these fix the ceiling
on what a practitioner may transfer.

\paragraph{What a practitioner should do differently} Three operational
consequences follow, each bounded by what was measured. Do not read a
reachability-gated relation's alarm as proof: because the reading time is part of
the specification, an alarm raised against a crawl-time snapshot---rather than
against live state---cannot be promoted to \texttt{violation-proven} on the
relation's word alone. If a readable, sound authorization oracle is in fact
available to the analyst, consult it directly and drop the relation: for the
relations studied it adds no adjudication information once that antecedent holds
(Proposition~\ref{lem:oracle-suff}, Section~\ref{sec:stratum}). And treat a
reported input-level specificity as a corpus statistic rather than a guarantee of
scenario-level resolution: it can sit near one while a whole class of legitimate
access is misreported (Finding~3). These are not a claim that metamorphic testing
is worthless; they locate its contribution where no oracle exists, and bound it
by what the reachability reading can see.

\paragraph{What a stronger claim would require} Moving from characterization to something
stronger needs, first, a population of real systems that \emph{contain} the
flaw, so detection can be measured---neither the four seeded classes of
Section~\ref{sec:multiclass} nor the single publicly fixed defect of
Section~\ref{sec:cvepair} supplies that population: the first are planted in an
emulator, the second is a real defect replayed under a controlled topology, and
neither is a sample of defects found in the field; and second, relations whose gate is not the absence of
evidence---checkable against a subject's entitlements without pretending that
entitlements are reachable facts. The transferable condition---we call it the
\emph{same-observable condition}---is the one that killed our own
mechanism: a component's necessity must be tested against a substitute that
consumes the \emph{same} observable quantity as the ground truth; a substitute
that consumes less does not establish incremental value. The \mbox{MST-wi} case
study is an instance of that condition, not its scope: a component's necessity is
not established unless a fair control that lacks the component but consumes the
same decision-relevant observables is empirically or formally shown to be
weaker.

% =====================================================================
% 声明章节（IEEE Access 目标；结构类无关，ieeeaccess / IEEEtran 均可编译）
%   IEEE Access 要求：Data availability / Competing interests / Funding /
%   Authors' contributions，以及生成式 AI 使用声明（另置 \section*）。
% 用 \section*/\subsection* 以保证类无关。
% 内容必须真实：未加任何待填占位符。
% =====================================================================

\section*{Declarations}

\subsection*{Data availability}

The raw records of the studies reported here are publicly available under the
version-pinned snapshot \texttt{artifact-v1} of the companion repository
\url{https://github.com/ubfbuz3/JISA} (re-anchored to the submission commit):
(1)~the scenario/configuration study
(\texttt{experiments/results/raw\_matrix.jsonl},
\texttt{ablation\_direct*.json}, \texttt{alt\_matrix.jsonl});
(2)~the real-system cell study across Gitea~\ResRsVersion{}, Gogs~\GgVersion{} and
GitLab~\GlVersion{} (\texttt{experiments/real\_system/results/*.json});
(3)~the execution of the published relations on the native engine
(\texttt{experiments/native\_engine/results/*.json});
(4)~the seeded positive control (\texttt{seeded.json}).
The analysis pipeline (\texttt{analysis/*.py}: \texttt{compute\_all.py} and the
ingest scripts) recomputes every number, table and figure in this paper from those
records; the stratified ablation of Section~\ref{sec:stratum} adds no new
experiment---it re-aggregates the same scenario/configuration records by whether
the introspection surface returned a reading, and is recomputed, together with
every macro it defines, by \texttt{ingest\_m2\_stratum.py};
\texttt{check\_macro\_use.py} asserts that every macro cited in the manuscript
has exactly one definition. Running \texttt{bash reproduce.sh} regenerates all
outputs and the PDF, and
\texttt{bash verify\_artifact.sh} asserts the SHA-256 integrity of every raw record
and byte-reproducible output. The manuscript \LaTeX{} source ships in the same
repository. A persistent archival copy with a citable DOI is deposited at
Zenodo: \url{https://doi.org/10.5281/zenodo.22859942}.

\subsection*{Competing interests}

The authors declare no competing interests, financial or otherwise.

\subsection*{Funding}

This research received no specific grant from any funding agency in the public,
commercial, or not-for-profit sectors.

\subsection*{Authors' contributions}

Both authors contributed to the work. Qiulong Lin and Nanlei Chen jointly
developed the conceptualisation and methodology, designed and implemented the
software and analysis pipeline, carried out the investigation and formal
analysis, and wrote the manuscript.

\subsection*{Ethics}

Not applicable. This study involved no human participants, no personal data, and
no animals. All experiments were performed against software systems that we
deployed and controlled locally.

\section*{Declaration of generative AI and AI-assisted technologies in the writing process}

During the preparation of this work the authors used Claude (Anthropic) to assist
with literature verification, the development of the analysis scripts, and
language editing. Where it contributed to analysis code, every generated or
modified script was reviewed, executed, and its output checked by the authors
against the raw records. After using this tool, the authors reviewed and edited
the content as needed and take full responsibility for the content of the
published article. All experimental computations and final reported values were
independently checked by the authors against the released scripts and raw
records.


\bibliographystyle{IEEEtran}
\bibliography{refs}

% ---------------------------------------------------------------------
%  Author biographies (IEEE Access requires one per author; they may also be
%  entered directly in the IEEE Author Portal). Facts confirmed by the authors
%  (2026-09-28). See author_bios_TODO.md.
% ---------------------------------------------------------------------
\begin{IEEEbiographynophoto}{Qiulong Lin}
Qiulong Lin received the M.S. degree from Sogang University, Seoul, Republic of Korea. He is with Yango University, Fuzhou, China, and is also affiliated with Sogang University, Seoul, Republic of Korea. His research interests include Web/API security, vulnerability discovery, and metamorphic testing.
\end{IEEEbiographynophoto}

\begin{IEEEbiographynophoto}{Nanlei Chen}
Nanlei Chen is currently pursuing the M.S. degree at Sogang University, Seoul, Republic of Korea. He is with Yango University, Fuzhou, China. His research interests include Web/API security and authorization testing. He is the corresponding author of this work.
\end{IEEEbiographynophoto}

% ieeeaccess.cls 强制要求正文末尾出现 \EOD（否则告警），并借它打印结尾标记。
\EOD

\end{document}
